% ---------------------------------------------------------------------------
% Chapitre 10 — Le produit StateValue et la stabilité référentielle
% Sources : notes/04-domains-lattices.md (§3.6–§3.9, §4.1, §4.5–§4.7,
% §5, §6.2–§6.8, §7, §8), src/domains/impls/{state_value,stability}.rs,
% src/domains/context.rs, src/domains/transfer/state_value.rs
% (eval_state_value, recompute_memo, as_arith, eval_field_access),
% src/rules/api/query.rs (StabilityVerdict), src/engine/written.rs
% (value_freshness) ; ADR-002, ADR-007, ADR-008, ADR-015, ADR-017, ADR-020.
% Sorties observées au commit e67b10a avec target/debug/reactant sur les
% fichiers /tmp/ch10/*.tsx ; valeurs abstraites relevées avec la sonde
% temporaire du dossier 04 (/tmp/dom04/target/debug/probe, hors dépôt).
% ---------------------------------------------------------------------------
\chapter{Le produit \code{StateValue} et la stabilité référentielle}
\label{chap:statevalue-stabilite}

\begin{objectifs}
  \item Comprendre pourquoi la valeur abstraite de \reactant{} est un \emph{produit} indexé par les sortes de JavaScript (ADR-015), et savoir calculer à la main son ordre, son join, son widening et sa concrétisation.
  \item Connaître le treillis \code{Stability} de l'ADR-017 : la trace de changement qu'il abstrait, les bornes may et must, les points \code{Stable}, \code{Versioned}, \code{VersionedTop}, \code{PerRender}, et savoir prouver la soundness de son join.
  \item Comprendre la double vue de l'état (le store range les valeurs écrites, la lecture rend une stabilité versionnée) et le traitement des littéraux, de \code{useRef}, \code{useMemo} et \code{useCallback}.
  \item Savoir quelle projection une règle doit employer (\code{to_stability}, \code{is_unstable_reference_only}, \code{is_unbounded}, \code{StabilityVerdict}), et pourquoi \code{PerRender} n'y a pas toujours le même sens.
  \item Connaître les contextes \code{AnalysisCtx}, \code{QueryContext}, \code{InterCtx} qui relient les domaines au moteur, et les défauts observés au commit \snapshotCommit.
\end{objectifs}

\prerequis{\cref{chap:treillis-simples} (les quatre treillis de base, le widening à seuils) ; \cref{chap:interpretation-abstraite} (\cref{def:treillis}, \cref{def:concretisation}, \cref{def:must-may}, \cref{def:soundness}) ; \cref{chap:react-modele} (\code{Object.is}, deps, bail-out) ; \cref{chap:ir} (\cref{def:slot}, \code{Expr::StateVal}).}

Le \cref{chap:treillis-simples} a décrit quatre petits domaines : intervalles,
booléens, ensembles bornés de chaînes, setters. Aucun ne suffit seul, parce
qu'une variable JavaScript n'a pas de type fixe : \code{useState(null)} rend
une valeur qui sera \code{null} au premier rendu et un nombre au suivant. Ce
chapitre assemble ces domaines en une seule valeur abstraite,
\code{StateValue}, et lui ajoute la pièce la plus propre à React : la
\emph{stabilité référentielle}, c'est-à-dire la réponse à la question
\og cette valeur sera-t-elle \code{Object.is}-égale à celle du rendu
précédent ?\fg{}.

Deux questions guident tout le chapitre. La première porte sur les
\emph{valeurs} : dans le compteur \code{setN(n + 1)} amorcé à \code{null},
l'analyseur voit-il que \code{n} grandit sans borne ? Elle mène au produit
ponctuel (\cref{sec:statevalue-stabilite-enum,sec:statevalue-stabilite-produit}).
La seconde porte sur l'\emph{identité} au fil des rendus : une dep qui lit un
état objet change-t-elle à chaque rendu ? Elle mène au treillis de stabilité
versionnée (\cref{sec:statevalue-stabilite-naif,sec:statevalue-stabilite-traces})
et à la double vue de l'état (\cref{sec:statevalue-stabilite-lecture}). Les
sections suivantes montrent comment les mémos, les règles et le moteur lisent
ces valeurs, puis recensent les limites observées.

% ===========================================================================
\section{Un enum plat perd \code{null} $\cup$ \code{number}}
\label{sec:statevalue-stabilite-enum}

\subsection{Le problème}

Voici un compteur sans garde, amorcé à \code{null} au lieu de \code{0} :

\begin{lstlisting}[language=TypeScript,caption={Un compteur nullable (\file{/tmp/ch10/nullable.tsx}, extrait)},label={lst:statevalue-stabilite-nullcounter}]
export function NullCounter() {
  const [n, setN] = useState(null);
  useEffect(() => {
    setN(n + 1);
  }, [n]);
  return <div>{n}</div>;
}
\end{lstlisting}

Concrètement, React monte le composant avec \code{n = null}, puis exécute
l'effet : \code{null + 1} vaut \code{1} en JavaScript (\code{ToNumber(null)}
vaut $0$). La dep \code{n} passe de \code{null} à \code{1}, l'effet se
ré-exécute, écrit \code{2}, puis \code{3}, et ainsi de suite : une boucle de
rendu infinie. Pour la détecter, l'analyseur doit voir la valeur de
\code{n} \emph{grandir}, c'est-à-dire garder l'intervalle numérique précis
même quand \code{null} fait aussi partie des valeurs possibles.

Le motif n'a rien d'exotique : \code{useState(null)} est la manière
idiomatique d'écrire \og pas encore chargé\fg{}. Un domaine qui ne sait pas
représenter \og \code{null} ou un nombre\fg{} laisse passer toute une famille
de boucles.

\subsection{Le domaine d'ADR-008 : une somme plate}

La première version du domaine des valeurs (\adr{8}, juin 2026) était une
somme : une valeur abstraite était \emph{exactement une} sorte JavaScript.

\begin{lstlisting}[language=Rust,caption={L'ancien \code{StateValue} (\file{docs/adr/ADR-008-value-domain.md}, élagué des commentaires)},label={lst:statevalue-stabilite-adr8}]
pub enum StateValue {
    Bottom,
    Null,
    Undefined,
    Number(Interval),
    Boolean(BoolVal),
    StrConst(Arc<BTreeSet<String>>),
    Str,
    Reference(Stability),
    Top,
}
\end{lstlisting}

Dans une somme, le join de deux sortes différentes n'a qu'un majorant
commun : $\Top$. L'ADR le dit en une ligne de son tableau :
\code{Null} $\join$ \code{Number(i)} $=$ \code{Top}
(\cref{fig:statevalue-stabilite-somme-produit}, à gauche). Sur
\code{NullCounter}, le premier tour écrit $[1,1]$ dans un slot qui vaut
\code{Null} ; le join donne $\Top$, et $\Top$ est un point fixe. Le store ne
s'élargit jamais, aucun widening n'a lieu, et le bras \og point fixe\fg{}
d'\regle{infinite-loop} (qui cherche un slot élargi,
\cref{chap:infinite-loop}) ne voit rien : un faux négatif (FN).

\begin{figure}[htbp]\centering
\begin{tikzpicture}[treillis,node distance=8mm and 4mm]
  % somme plate
  \node (top) at (0,2.2) {$\Top$};
  \node[font=\scriptsize] (num) at (-3.3,1.1) {\code{Number}};
  \node[font=\scriptsize] (bool) at (-2,1.1) {\code{Boolean}};
  \node[font=\scriptsize] (str) at (-0.7,1.1) {\code{Str}\dots};
  \node[font=\scriptsize] (ref) at (0.7,1.1) {\code{Reference}};
  \node[font=\scriptsize,text=ra-blue] (null) at (2,1.1) {\code{Null}};
  \node[font=\scriptsize] (undef) at (3.3,1.1) {\code{Undefined}};
  \node (bot) at (0,0) {$\Bot$};
  \foreach \x in {num,bool,str,ref,null,undef} {\draw (top) -- (\x.north); \draw (\x.south) -- (bot);}
  \node[etiquette,align=center] at (0,-0.8) {somme (ADR-008) :\\ \code{Null} $\join$ \code{Number([1,1])} $=\Top$};
  % produit
  \begin{scope}[xshift=7.2cm]
    \node[cfgbloc] (a) at (-1.2,1.9) {num : $\Bot$\\ null : vrai};
    \node[cfgbloc] (b) at (1.2,1.9) {num : $[1,1]$\\ null : faux};
    \node[cfgbloc] (j) at (0,0.3) {num : $[1,1]$\\ null : vrai};
    \draw[fleche] (a) -- (j); \draw[fleche] (b) -- (j);
    \node[etiquette,align=center] at (0,-0.8) {produit (ADR-015) :\\ \code{number[1, 1]|null}};
  \end{scope}
\end{tikzpicture}
\caption{Le même join dans les deux domaines. À gauche, la somme plate
d'ADR-008 : deux sortes différentes n'ont que $\Top$ pour majorant commun
(les sous-treillis de chaque sorte sont repliés en un nœud). À droite, le
produit d'ADR-015 : chaque sorte garde sa composante, le join se fait
composante par composante (seules deux composantes sont montrées).}
\label{fig:statevalue-stabilite-somme-produit}
\end{figure}

\subsection{Trois rustines, et un FN qui résiste}

L'ADR-008 a compensé la perte par trois mécanismes empilés
(\adr{15}, section \og Context\fg{}) :
\begin{enumerate}
  \item un \code{TypedStateStore} à cinq sous-stores (nombres, booléens,
    chaînes, références, inconnus), choisis label par label ;
  \item une inférence \code{infer_state_type} qui devinait la sorte d'un
    slot à partir de son initialiseur ;
  \item la lecture de l'annotation TypeScript \code{useState<number>(null)} :
    le lowering la capturait dans \code{HookEntry::State.type_hint}, et le
    point fixe remplaçait l'amorçage \code{Null} par \code{Number([0,0])}.
\end{enumerate}
Même avec les trois, \code{useState(null)} \emph{sans} annotation restait un
FN documenté, et un code JavaScript pur qui mélange nombres et chaînes
(\code{let a = 10; if (c) a = "s";}) perdait toute précision. Les rustines
traitaient le symptôme (une sorte par slot) sans toucher à la cause (une
sorte par valeur).

\begin{decision}[ADR-008 --- domaine des valeurs en somme plate]
\textbf{Décidait} : un \code{StateValue} en \code{enum} plat, tout join
inter-sortes vaut $\Top$ ; précision regagnée par un store typé et par les
annotations TypeScript. \textbf{Statut} : remplacée par \adr{15}
(14 juillet 2026). \textbf{Ce qui survit} : les sous-domaines
\code{Interval}, \code{BoolVal}, \code{StrConst} (seuil 4), la machinerie de
widening et de raffinement, et le signal \og le store s'élargit, donc boucle
possible\fg{}. \textbf{À ne pas croire} : la ligne du tableau d'ADR-008
\og Boolean : oscillation true$\leftrightarrow$false détectée\fg{} ne
correspond pas au comportement actuel (\cref{sec:statevalue-stabilite-limites}).
\end{decision}

% ===========================================================================
\section{Sortes disjointes : le produit \code{StateValue}}
\label{sec:statevalue-stabilite-produit}

\subsection{L'observation}

Une valeur JavaScript n'est jamais à la fois un nombre et une chaîne, ni à
la fois \code{null} et un objet : les sortes sont \emph{disjointes}. Or ce
que l'on veut abstraire, c'est un \emph{ensemble} de valeurs possibles
$V$. Un tel ensemble se découpe de manière unique selon les sortes :
\[
  V \;=\; (V \cap \mathit{Num}) \uplus (V \cap \mathit{Bool}) \uplus
  (V \cap \mathit{Str}) \uplus \dots
\]
Chaque morceau s'abstrait dans le treillis de sa sorte, indépendamment des
autres. Le domaine idéal (la \terme{complétion disjonctive} de la somme, qui
représente des disjonctions de valeurs abstraites) dégénère donc, pour une
somme de sortes disjointes, en un \terme{produit ponctuel} : une composante
par sorte, chacune à $\Bot$ quand la sorte est impossible. C'est
l'observation d'ADR-015, et elle tient en une phrase, conformément au
principe du paragraphe unique de \file{CLAUDE.md}.

\begin{remarque}
Dans une somme non disjointe (par exemple \og entier\fg{} et \og entier
pair\fg{}), une même valeur pourrait relever de deux composantes, et un
produit perdrait la corrélation entre elles. C'est précisément la
disjonction des sortes qui rend le produit exact, sans opérateur de
réduction entre composantes.
\end{remarque}

\subsection{Le code}

\rustsnippet[label={lst:statevalue-stabilite-struct}]{src/domains/impls/state_value.rs}{16}{44}{le produit ponctuel des sortes JavaScript}

Les huit composantes, champ par champ :
\begin{itemize}
  \item \code{num} : un \code{Interval} (\cref{sec:treillis-simples-intervalles}) ;
    $\Bot$ signifie \og ne peut pas être un nombre\fg{}.
  \item \code{boolean} : le treillis plat \code{BoolVal}
    (\cref{sec:treillis-simples-plats}).
  \item \code{str} : l'ensemble borné \code{StrConst}, élargi à $\Top$ au-delà de
    quatre constantes (\cref{sec:treillis-simples-strconst}).
  \item \code{reference} : la sorte \og objet, tableau ou fonction\fg{}. Sa
    composante n'est pas un ensemble de valeurs mais une
    \emph{stabilité} (\cref{sec:statevalue-stabilite-traces}) ; $\Bot$
    signifie \og ne peut pas être une référence\fg{}.
  \item \code{null}, \code{undef} : deux booléens Rust, \code{true} voulant
    dire \og cette valeur est possible\fg{}. L'ordre de \code{bool}
    (\code{false < true}) est exactement le treillis à deux points
    $\Bot \lleq \Top$.
  \item \code{setter} : le treillis plat \code{SetterVal} des setters
    d'état identifiés par \code{(ComponentId, HookLabel)}.
  \item \code{other} : le \terme{résidu}, \og tout le reste\fg{} (symbol,
    bigint, et toute sorte non modélisée).
\end{itemize}

\begin{definition}{Le produit \code{StateValue}}{statevalue}
  Une valeur abstraite est un octuplet
  \[
  v = (v.\mathit{num}, v.\mathit{boolean}, v.\mathit{str},
  v.\mathit{reference}, v.\mathit{null}, v.\mathit{undef},
  v.\mathit{setter}, v.\mathit{other}),
  \]
  élément du produit des huit treillis ci-dessus. L'ordre, le join, le meet et le widening sont
  \emph{ponctuels} (composante par composante). La concrétisation est
  l'union disjointe des concrétisations des composantes :
  \[
  \begin{array}{rl}
    \gam(v) \;=& \gam_{\mathit{num}}(v.\mathit{num}) \;\cup\;
                 \gam_{\mathit{bool}}(v.\mathit{boolean}) \;\cup\;
                 \gam_{\mathit{str}}(v.\mathit{str})\\[2pt]
       \cup& \{\, o \text{ objet, tableau ou fonction} \mid
                 \text{la trace de changement de } o \in
                 \gam_{\Stab}(v.\mathit{reference}) \,\}\\[2pt]
       \cup& \{\code{null} \mid v.\mathit{null}\} \;\cup\;
                 \{\code{undefined} \mid v.\mathit{undef}\} \;\cup\;
                 \gam_{\mathit{setter}}(v.\mathit{setter})\\[2pt]
       \cup& \{\, \text{toute autre valeur (symbol, bigint, \dots)} \mid
                 v.\mathit{other} \,\}
  \end{array}
  \]
  où $\gam_{\Stab}$ est la concrétisation de la stabilité
  (\cref{def:stabilite}).
\end{definition}

\begin{exemple}{Lire une valeur du produit}{statevalue-stabilite-lire}
  La valeur affichée \code{number[1, inf]|null} a
  $\mathit{num} = [1, +\infty]$, $\mathit{null} = \code{true}$, et toutes
  les autres composantes à $\Bot$. Sa concrétisation est
  $\{x \in \mathbb{R} \mid x \ge 1\} \cup \{\code{null}\}$ (avec le bit
  \code{is_int}, les seuls entiers). Elle ne contient ni \code{undefined}, ni
  chaîne, ni objet. L'affichage (\code{Debug}, écrit à la main,
  \src{src/domains/impls/state_value.rs}{422}{476}) juxtapose les sortes
  peuplées séparées par \texttt{|}.
\end{exemple}

\begin{piege}
  La composante \code{reference} ne ressemble pas aux autres. Pour
  \code{num} ou \code{str}, $\gam$ décrit des \emph{valeurs} ; pour
  \code{reference}, elle décrit la \emph{stabilité} des références possibles,
  sans rien dire de leur contenu. Le contenu des objets relève du tas par
  site d'allocation (\cref{chap:transfert-stores}). Une \code{StateValue}
  \code{ref(Stable)} dit \og une référence qui ne change pas d'un rendu à
  l'autre\fg{}, pas \og l'objet \code{{a: 1}}\fg{}.
\end{piege}

Et \code{NaN} ? Aucune composante n'en parle. Toute opération qui pourrait le
produire (division, reste par un diviseur qui peut être nul, puissance,
\code{undefined} en arithmétique) rend $\Top$, dont le résidu \code{other}
est vrai : c'est par là que \code{NaN} est couvert
(\cref{sec:treillis-simples-arith}). Le code le dit explicitement pour
\code{\%} et \code{**} : \og where the result may be \code{NaN}, which no
interval holds, they are $\Top$ like \code{/}\fg{}
(\srcline{src/domains/transfer/state_value.rs}{781}).

\subsection{$\Bot$, $\Top$ et les constructeurs}

\rustsnippet[label={lst:statevalue-stabilite-bottom}]{src/domains/impls/state_value.rs}{190}{206}{$\Bot$ : toutes les composantes à $\Bot$}

\code{bottom_value} est une \code{const fn} pour pouvoir servir dans la
syntaxe de mise à jour de structure (\code{..Self::bottom_value()}) de
chaque constructeur. L'intervalle vide y est écrit à la main avec ses
bornes sentinelles $[+\infty, -\infty]$ : tout prédicat qui lit
\code{num.lo} ou \code{num.hi} doit donc d'abord tester \code{is_bottom()}
(\cref{sec:statevalue-stabilite-projections}). $\Top$ met toutes les
composantes à leur $\Top$, \emph{y compris} le résidu :

\rustsnippet[label={lst:statevalue-stabilite-top}]{src/domains/impls/state_value.rs}{512}{523}{$\Top$ : toutes les composantes à $\Top$, résidu compris}

Les constructeurs \code{number}, \code{boolean}, \code{str_set},
\code{str_singleton}, \code{str_top}, \code{reference}, \code{null},
\code{undefined} et \code{component_setter}
(\src{src/domains/impls/state_value.rs}{125}{188}) peuplent une seule
composante et laissent les autres à $\Bot$ :

\rustsnippet{src/domains/impls/state_value.rs}{125}{130}{la forme commune des constructeurs}

Deux exceptions : \code{str_set} passe par \code{StrConst::from_set}, donc
un ensemble vide donne $\Bot$ et plus de quatre chaînes donnent
\code{str: Top} ; \code{component_setter(c, l)} pose
\code{SetterVal::One(c, l)}.

\subsection{\code{KindMask} : une seule énumération des sortes}

Beaucoup de questions sur une valeur portent sur \emph{quelles} sortes sont
peuplées : la valeur est-elle $\Bot$ ? n'est-elle qu'un setter ? combien de
sortes mêle-t-elle ? Si chaque prédicat ré-énumérait les huit champs,
l'ajout d'un neuvième champ serait silencieusement oublié quelque part, et
un prédicat comme \og la valeur n'est qu'une référence fraîche\fg{}
deviendrait vrai à tort : un faux négatif latent. Le code centralise donc
l'énumération dans un masque de bits.

\rustsnippet[label={lst:statevalue-stabilite-kindmask}]{src/domains/impls/state_value.rs}{46}{78}{le masque des sortes peuplées}

\rustsnippet[label={lst:statevalue-stabilite-populated}]{src/domains/impls/state_value.rs}{80}{121}{l'unique énumération des sortes}

\begin{invariant}{Une seule énumération des sortes}{statevalue-stabilite-kindmask}
  Tout prédicat de la forme \og quelles sortes sont peuplées\fg{} se dérive
  de \code{populated_kinds}, dont la déstructuration n'a pas de \code{..}.
  Ajouter une composante à \code{StateValue} est donc une erreur de
  compilation dans \code{populated_kinds} (et dans \code{is_top_value}, qui
  utilise la même astuce pour la saturation à $\Top$), pas un oubli
  silencieux.
\end{invariant}

Les consommateurs du masque sont \code{is_bottom_value} (\code{is_empty}),
\code{typeof_name} (une seule sorte), \code{as_setter} et
\code{is_unstable_reference_only} (\code{only}), et \code{to_stability}
(\code{count() >= 2}). La protection vient d'ADR-020 (\adr{20}, décision D4,
commit \code{c887c52}).

\rustsnippet[label={lst:statevalue-stabilite-as-setter}]{src/domains/impls/state_value.rs}{270}{287}{deux prédicats dérivés du masque}

\begin{piege}
  \code{only(bit)} est \emph{vrai} sur $\Bot$ : aucun bit n'est peuplé, donc
  tous les bits peuplés sont bien dans \code{{bit}}. Chaque consommateur
  doit écarter $\Bot$ lui-même s'il le faut. \code{is_unstable_reference_only}
  le fait en exigeant en plus \code{reference == PerRender} ;
  \code{as_setter} le fait parce que \code{SetterVal::as_one} rend
  \code{None} sur $\Bot$ et $\Top$.
\end{piege}

\begin{piege}
  La protection est partielle. Les fonctions de transfert \code{as_arith},
  \code{as_str_only} (\src{src/domains/transfer/state_value.rs}{708}{750})
  et la négation \code{!} (\src{src/domains/transfer/state_value.rs}{947}{965}) ré-énumèrent les
  composantes à la main, sans passer par \code{KindMask} : une composante
  ajoutée ne les ferait pas échouer à la compilation (constat de lecture,
  dossier 04 §3.7).
\end{piege}

\code{typeof_name} illustre la sémantique JavaScript encodée dans le
masque : \code{typeof null} vaut \code{"object"}, un setter est une
\code{"function"}, et la composante \code{reference} ne donne \emph{pas} de
réponse, parce qu'elle couvre à la fois les objets (\code{"object"}) et les
fonctions (\code{"function"})
(\src{src/domains/impls/state_value.rs}{240}{268}).

\subsection{L'arithmétique : \code{as_arith} et \code{ToNumber(null)}}

Le produit ne sert que si l'évaluateur sait calculer \code{n + 1} quand
\code{n} vaut \code{number[1, 1]|null}. La règle est celle de JavaScript :
\code{ToNumber(null)} vaut $0$. L'évaluateur demande donc une \emph{vue
numérique} de chaque opérande :

\rustsnippet[label={lst:statevalue-stabilite-as-arith}]{src/domains/transfer/state_value.rs}{708}{732}{la vue numérique d'un opérande}

\code{as_arith} rend un intervalle si les seules sortes peuplées sont
\og nombre\fg{} et \og \code{null}\fg{} ; la contribution de \code{null} est
le point $[0,0]$, joint à l'intervalle par \code{hull}. Toute autre sorte
(\code{undefined}, qui donnerait \code{NaN} ; booléen ; chaîne ; référence ;
setter ; résidu) rend \code{None}, et l'opération arithmétique vaut alors
$\Top$ (\cref{sec:treillis-simples-arith}).

\begin{proposition}{Soundness de la vue numérique}{statevalue-stabilite-as-arith}
  Si \code{as_arith(v)} $= \mathit{Some}(I)$, alors pour toute valeur
  concrète $x \in \gam(v)$, $\code{ToNumber}(x) \in \gam(I)$.
\end{proposition}
\begin{proof}
  $\mathit{Some}$ n'est rendu que si toutes les composantes autres que
  \code{num} et \code{null} sont à $\Bot$ ; donc $\gam(v) \subseteq
  \gam(v.\mathit{num}) \cup \{\code{null}\}$. Un nombre est son propre
  \code{ToNumber} et appartient à $\gam(v.\mathit{num}) \subseteq
  \gam(I)$ ; \code{null} donne $0 \in \gam([0,0]) \subseteq \gam(I)$
  lorsque $v.\mathit{null}$ est vrai, puisque \code{hull} est un majorant.
\end{proof}

Le commentaire \og a $\Bot$ interval stays \code{Some(}$\Bot$\code{)}\fg{} n'est pas un
détail. Sur une branche morte, le raffinement de garde met \code{num} à
$\Bot$ ; l'addition rend alors $\Bot$, et l'écriture de $\Bot$ dans un
store se joint sans effet. Si \code{as_arith} rendait \code{None}, la même
écriture vaudrait $\Top$ et polluerait le slot.

\begin{exemple}{\code{NullCounter}, tour par tour}{statevalue-stabilite-nullcounter}
  Amorçage : $n \mapsto \{\code{null}\}$. Premier tour : l'effet évalue
  \code{n + 1} ; \code{as_arith({null})} $= [0,0]$, donc l'écriture vaut
  $[1,1]$ ; le store fait le join ponctuel et contient
  \code{number[1, 1]|null}. Deuxième tour : \code{as_arith} rend
  $[0,1]$, l'écriture $[1,2]$, le store \code{number[1, 2]|null}. Au
  troisième tour, la boucle externe applique \code{widen_to}
  (\cref{sec:treillis-simples-widening}), qui n'agit que sur \code{num} ;
  la borne haute saute à $+\infty$, \code{null} reste vrai. La sonde relève
  la valeur convergée :
\begin{sortie}
-- NullCounter
   state[0] = number[1, inf]|null   widened=true  num=Interval { lo: 1.0, hi: inf, is_int: true }
\end{sortie}
  et la ligne de commande signale la boucle :
\begin{sortie}
$ reactant check /tmp/ch10/nullable.tsx --trace
  NullCounter  (2 hooks)  /tmp/ch10/nullable.tsx
    warn   infinite-loop  [hook:0]  (line 5:2)  this effect keeps pushing state `n` (its deps do not provably gate it, so the effect can re-run every render) to new values on every run. Potential infinite render loop
       → state `n` is written here [hook:1] (line 5:2)
       → the abstract value of state `n` kept growing and was widened at iteration 3
\end{sortie}
  Le FN d'ADR-008 est devenu une détection, épinglée par le test
  \code{tests/narrowing.rs::null_init_without_hint_unbounded_is_flagged}.
\end{exemple}

La \og sonde\fg{} citée ici et dans la suite du chapitre est un petit
programme temporaire, hors du dépôt, qui appelle \code{lower_program} puis
\code{analyze_component} sur un fichier et imprime les valeurs convergées
des stores (\code{state[l]} pour le slot de label $l$, \code{memo[l]} pour
un mémo), avec le drapeau \code{widened} (le slot est-il dans
\code{widen_trace} ?). Les identités de composant y apparaissent comme
\code{ComponentId(4294967295)}, c'est-à-dire \code{ComponentId::SYNTHETIC}
d'une analyse isolée.

Deux variantes du même fichier montrent les bords du mécanisme. Remplacer
\code{useState(null)} par \code{useState(undefined)} (composant
\code{UndefCounter}) fait échouer \code{as_arith} : \code{undefined + 1}
vaut \code{NaN}, que l'intervalle ne sait pas représenter, donc l'écriture
vaut $\Top$ et le store aussi (la sonde affiche \code{state[0] =} $\Top$). L'analyseur
émet alors un \Warning{} plus vague :
\begin{sortie}
  UndefCounter  (2 hooks)  /tmp/ch10/nullable.tsx
    warn   infinite-loop  [hook:0]  (line 21:2)  this effect may store a fresh reference into state `n` which its deps react to: possible infinite render loop
       → a fresh value is written to state `n` here [hook:1] (line 22:4)
\end{sortie}
C'est un faux positif (FP) toléré : concrètement, \code{n} passe de
\code{undefined} à \code{NaN}, puis \code{NaN + 1} vaut \code{NaN}, et
\code{Object.is(NaN, NaN)} est vrai, donc React abandonne le rendu suivant.
La valeur $\Top$ a une composante \code{reference} à \code{Unknown}, que le
graphe de churn lit comme \og peut-être fraîche\fg{}
(\cref{sec:statevalue-stabilite-lecture}). Le motif \og chargement
unique\fg{} \code{if (!user) setUser({ name: "guest" })}, lui, converge
sans bruit (le composant \code{FetchOnce} est propre) : la branche
\code{!user} applique \code{narrow_falsy}, qui tue la composante
\code{reference} (\cref{sec:treillis-simples-gardes}), et l'écriture n'a
lieu qu'une fois ; le store se stabilise sur \code{ref(PerRender)|null}.

\subsection{Ordre, join, meet et widening ponctuels}

\rustsnippet[label={lst:statevalue-stabilite-ord}]{src/domains/impls/state_value.rs}{479}{501}{l'ordre produit}

\code{merge_ord} combine les comparaisons des composantes : \code{Equal} est
neutre, deux \code{Less} (ou deux \code{Greater}) restent dans leur sens,
et deux sens opposés, ou une composante incomparable (\code{None}, propagé
par l'opérateur \code{?}), rendent tout le produit incomparable. On a donc
$a \lleq b$ si et seulement si chaque composante de $a$ est sous celle de
$b$ : c'est l'ordre produit (test \code{mixed_direction_slots_are_incomparable}).

\rustsnippet[label={lst:statevalue-stabilite-join}]{src/domains/impls/state_value.rs}{631}{676}{join, meet, widen et widen\_to ponctuels}

Trois remarques de lecture :
\begin{itemize}
  \item le join de \code{num} est \code{hull}, celui des trois booléens est
    le \og ou\fg{}, celui des autres composantes est le join de leur
    treillis ;
  \item le widening n'élargit que les composantes de hauteur infinie ou
    grande : \code{num} (\code{Interval::widen}) et \code{str}
    (\code{StrConst::widen}) ; les autres se contentent du join ;
  \item \code{widen_to} ne passe les seuils qu'à \code{num} et hérite du
    reste par la syntaxe \code{..self.widen(other)}.
\end{itemize}

\begin{theoreme}{Soundness du join et du widening du produit}{statevalue-stabilite-join}
  Pour toutes valeurs $a, b$ : $\gam(a) \cup \gam(b) \subseteq \gam(a
  \join b)$ et $\gam(a) \cup \gam(b) \subseteq \gam(a \widen b)$. De plus,
  toute suite $x_0, x_1 = x_0 \widen y_1, x_2 = x_1 \widen y_2, \dots$ est
  stationnaire.
\end{theoreme}
\begin{proof}
  $\gam$ est une union, composante par composante, de concrétisations
  disjointes (\cref{def:statevalue}). Il suffit donc que chaque composante
  soit sound : $\gam_k(a_k) \cup \gam_k(b_k) \subseteq \gam_k(a_k \join
  b_k)$. C'est vrai pour \code{hull}, \code{BoolVal}, \code{StrConst},
  \code{SetterVal} (\cref{sec:treillis-simples-soundness}), pour les trois
  booléens (le \og ou\fg{} est l'union de $\emptyset$ et du singleton) et
  pour \code{Stability} (\cref{thm:statevalue-stabilite-join-stab}). Le
  widening de chaque composante majore son join, d'où la seconde inclusion.
  Terminaison : une suite de widenings est stationnaire composante par
  composante, parce que \code{num} a un vrai widening (le
  \cref{chap:treillis-simples} le montre), que \code{str} saute à $\Top$
  au-delà de quatre constantes, et que les six autres composantes sont de
  hauteur finie (au plus 8 pour \code{Stability}). Un produit fini de suites
  stationnaires est stationnaire.
\end{proof}

Le meet est lui aussi ponctuel ; il n'est employé par aucun code de
production (\cref{sec:treillis-simples-contrat}), ce qui compte pour la
stabilité (\cref{sec:statevalue-stabilite-traces}).

\begin{decision}[ADR-015 --- produit ponctuel sur les sortes disjointes]
\textbf{Décide} : \code{StateValue} devient le produit de huit composantes ;
\code{join}, \code{meet}, \code{widen}, \code{widen_to} et
\code{partial_cmp} sont ponctuels ; \code{ComponentSetter} devient une
composante (\code{SetterVal}) ; l'arithmétique applique
\code{ToNumber(null) = 0}. \textbf{Supprime} : \code{TypedStateStore},
\code{StateType}, \code{infer_state_type}, \code{HookEntry::State.type_hint} ;
\code{useState<T>(...)} devient décoratif. \textbf{Refuse} : un opérateur de
réduction ou un \og query pool\fg{} à la MOPSA \emph{à l'intérieur} du
produit, puisque des sortes disjointes n'ont aucune information à
échanger ; un futur produit réduit (valeur $\times$ relationnel) se fera au
niveau de l'analyse, avec \code{QueryContext} comme canal de lecture
(\cref{sec:statevalue-stabilite-contextes}). \textbf{Épingle} :
\code{null_init_without_hint_unbounded_is_flagged},
\code{nullable_fetch_pattern_no_false_positive}.
\end{decision}

\subsection{\code{from_init} : un nom trompeur}

\rustsnippet[label={lst:statevalue-stabilite-from-init}]{src/domains/impls/state_value.rs}{289}{306}{la valeur syntaxique d'un initialiseur}

\begin{piege}
  Malgré son nom et sa documentation, \code{from_init} n'est \emph{pas} la
  fonction qui amorce les slots \code{useState}. L'amorçage évalue
  l'initialiseur avec l'évaluateur complet (\code{transfer.eval_expr}), et
  exécute le corps d'un initialiseur paresseux \code{useState(() => expr)}
  au lieu d'abstraire la fermeture en référence fraîche
  (\src{src/engine/fixpoint.rs}{314}{353}, \cref{chap:point-fixe-composant}).
  Le seul appelant hors tests de \code{from_init} est
  \code{returns_value} (\src{src/engine/written.rs}{108}{121}), qui calcule
  sans environnement la valeur qu'un updater \code{setX(prev => ...)} rend
  sur ses \code{return}. Pour cet usage purement syntaxique, un
  \code{FnLit} rendu vaut bien \code{PerRender} : un updater qui rend une
  fonction stocke une nouvelle fonction.
\end{piege}

% ===========================================================================
\section{La stabilité référentielle : un premier treillis, et son impasse}
\label{sec:statevalue-stabilite-naif}

\subsection{Ce que React compare}

\begin{react}[Identité des valeurs d'un rendu à l'autre]
  React compare chaque entrée d'un tableau de deps à celle du rendu
  précédent avec \code{Object.is} (comme \code{===}, sauf
  \code{Object.is(NaN, NaN)} vrai et \code{Object.is(0, -0)} faux). Un effet,
  un \code{useMemo} ou un \code{useCallback} est ré-exécuté si \emph{au moins
  une} dep diffère. Or :
  \begin{itemize}
    \item \code{{}}, \code{[]}, \code{() => {}}, \code{new X()},
      \code{arr.map(f)} écrits dans le corps du composant allouent une
      nouvelle référence à chaque évaluation, donc à chaque rendu ;
    \item un état garde son identité d'un rendu au suivant tant que son
      setter n'est pas appelé avec une valeur différente ;
    \item la fonction setter d'un \code{useState} et le conteneur rendu par
      \code{useRef} sont les mêmes à chaque rendu (garantie de React) ;
    \item \code{useMemo} et \code{useCallback} rendent la même référence tant
      que leurs deps sont \code{Object.is}-égales.
  \end{itemize}
  Une primitive, elle, est comparée par valeur : \code{1} est toujours
  \code{Object.is}-égal à \code{1}, quelle que soit la façon dont on l'a
  calculé (\cref{chap:react-modele}).
\end{react}

La \terme{stabilité référentielle} d'une expression est la réponse
abstraite à la question \og sa valeur est-elle la même, au sens de
\code{Object.is}, que celle du rendu précédent ?\fg{}. C'est le domaine
qu'\adr{2} plaçait au centre de l'analyse, dès la première semaine du
projet.

\subsection{Le treillis d'ADR-002}

L'ADR-002 proposait un treillis plat à quatre points, de la forme étudiée au
\cref{sec:treillis-simples-plats} : $\Bot \lleq \{\code{Stable},
\code{Unstable}\} \lleq \code{Unknown}$, avec $\code{Stable} \join
\code{Unstable} = \code{Unknown}$. Une table de transfert statique en
donnait la valeur : un littéral primitif, un setter et \code{useRef()}
étaient \code{Stable} ; un littéral objet, tableau ou fonction était
\code{Unstable} ; la valeur d'un \code{useState} était le join des
arguments de ses setters ; \code{useMemo} et \code{useCallback}, le join de
leurs deps.

\begin{decision}[ADR-002 --- treillis de stabilité et trois stores]
\textbf{Décidait} : le domaine central est la stabilité (\code{Object.is}),
treillis à quatre points ; trois stores séparés (état soumis au point fixe,
mémos fonctions de leurs deps, refs triviales) parce que les trois familles
de hooks n'ont pas la même sémantique vis-à-vis du cycle de rendu ; widening
après un seuil configurable. \textbf{Alternative refusée} : un store unifié,
qui aurait forcé tous les hooks dans le même point fixe. \textbf{Ce qui a
changé depuis}, sans que le fichier porte la mention \og Superseded\fg{} :
\code{Unstable} est devenu \code{PerRender} et le treillis a gagné
\code{Versioned} et \code{VersionedTop} (\adr{17}) ; il n'existe pas de
\code{RefStore} (\code{useRef} rend \code{reference(Stable)}) ; les fichiers
\file{src/domains/stability.rs} et \file{src/domains/product.rs} annoncés
n'existent pas ; le seuil de widening par défaut est 3 et non 2.
\end{decision}

\subsection{Le faux positif F5}

Le premier banc d'essai sur un vrai corpus (\og memos\fg{}, 15 juillet 2026)
a laissé cinq \Warning{} \regle{always-unstable-deps} de la même forme, un
fournisseur de contexte qui tient un état objet :

\begin{lstlisting}[language=TypeScript,caption={Un état objet en dep (\file{/tmp/ch10/stability.tsx}, extrait)},label={lst:statevalue-stabilite-objstate}]
export function ObjectStateDep() {
  const [ctx, setCtx] = useState({ locale: "en" });
  useEffect(() => { console.log(ctx); }, [ctx]);
  return <button onClick={() => setCtx({ locale: "fr" })}>fr</button>;
}
\end{lstlisting}

L'initialiseur est un littéral objet, donc \code{Unstable} ; le handler écrit
un autre littéral, \code{Unstable} aussi ; le store du slot contient
\code{Unstable}, et chaque lecture de \code{ctx} rendait cette valeur. La
règle concluait \og \code{ctx} est une nouvelle référence à chaque
rendu\fg{}. C'est faux : \code{ctx} ne change que lorsqu'on clique. Le
domaine confondait deux faits différents :
\begin{itemize}
  \item la \emph{fraîcheur à l'allocation}, un fait sur un événement : la
    valeur \emph{écrite} par \code{setCtx({...})} est un objet neuf ;
  \item la \emph{fraîcheur par rendu}, un fait entre deux rendus : la valeur
    \emph{lue} par \code{ctx} change à chaque rendu.
\end{itemize}
La première est vraie, la seconde non, et l'ancien \code{Unstable} disait
les deux à la fois.

\subsection{Le faux négatif caché dessous}

Corriger le FP \og naïvement\fg{}, en rendant \og plutôt stable\fg{} toute
lecture d'état, aurait ouvert un faux négatif réel. L'ADR-017 le montre
avec ce composant :

\begin{lstlisting}[language=TypeScript,caption={Une boucle de churn d'objet (\file{/tmp/ch10/stability.tsx}, extrait)},label={lst:statevalue-stabilite-objchurn}]
export function ObjChurn() {
  const [obj, setObj] = useState({ a: 1 });
  useEffect(() => { setObj({ ...obj, b: 2 }); }, [obj]);
  return <div />;
}
\end{lstlisting}

L'effet écrit à chaque exécution un objet neuf dans \code{obj} ; la dep
\code{obj} change donc, l'effet se ré-exécute, et ainsi de suite : une vraie
boucle infinie. Or le bras point fixe d'\regle{infinite-loop} ne la voyait
pas, parce que le store ne grandit pas : $\code{Unstable} \join
\code{Unstable} = \code{Unstable}$, rien à élargir, contrairement à un
intervalle. Le \emph{seul} signal sur ce bug était le même \Warning{}
\regle{always-unstable-deps} qui était un FP sur \code{ObjectStateDep}.
Faire taire l'un sans remplacer le signal faisait disparaître l'autre, et
un FN est interdit (\cref{def:soundness}). La correction F5 devait donc être
double : scinder le treillis, \emph{et} ajouter une détection du churn qui
ne repose pas sur la divergence du store.

% ===========================================================================
\section{Traces de changement, bornes may et must}
\label{sec:statevalue-stabilite-traces}

\subsection{L'objet concret : la trace de changement}

Pour dire proprement ce que \og stable\fg{} veut dire, il faut fixer ce que
l'on abstrait. Ce n'est pas une valeur, mais le comportement d'une
expression au fil des rendus.

\begin{definition}{Trace de changement}{statevalue-stabilite-trace}
  Soit une expression d'un composant (une dep, une variable) évaluée aux
  rendus successifs $r_0, r_1, \dots, r_n$, avec les valeurs $v_0, \dots,
  v_n$. Sa \terme{trace de changement} est l'ensemble des rendus où
  \code{Object.is} échoue :
  \[
    C \;=\; \{\, i \in \{1,\dots,n\} \mid \neg\,\code{Object.is}(v_i, v_{i-1}) \,\}.
  \]
  Pour un ensemble $S$ de slots d'état, on note $\mathit{Set}(S)$ l'ensemble
  des rendus $r_i$ ($i \ge 1$) précédés, depuis $r_{i-1}$, d'au moins un
  appel au setter d'un slot de $S$.
\end{definition}

La trace de changement est exactement ce que React observe : les deps d'un
effet, l'invalidation d'un mémo et la propagation d'un contexte ne dépendent
que d'elle.

\begin{exemple}{Quatre traces}{statevalue-stabilite-traces}
  \begin{itemize}
    \item \code{const r = useRef()} : $C = \emptyset$ (même conteneur à
      chaque rendu).
    \item \code{const o = { a: 1 }} dans le corps : $C = \{1, \dots, n\}$
      (un objet neuf à chaque rendu).
    \item \code{ctx} lu depuis \code{useState({ locale: "en" })} dans
      \code{ObjectStateDep} : $C \subseteq \mathit{Set}(\{\code{ctx}\})$ ; il
      ne change qu'aux clics.
    \item \code{const v = flag ? { a: 2 } : obj}, avec \code{obj} un état :
      tantôt neuf à chaque rendu, tantôt versionné ; on ne peut rien borner.
  \end{itemize}
\end{exemple}

\subsection{Deux bornes, deux usages}

Sur l'ensemble $C$, deux informations utiles coexistent. Elles ne sont pas
opposées : ce sont deux \emph{bornes} de nature différente, au sens de la
\cref{def:must-may}.
\begin{itemize}
  \item Une borne \terme{may} (sur-approximation) : \og $C \subseteq X$\fg{},
    par exemple \og ne change \emph{qu'aux} sets de $S$\fg{}. C'est un fait
    de sûreté. Les règles s'en servent pour \emph{se taire} à bon droit : une
    dep versionnée ne ré-exécute pas l'effet à chaque rendu.
  \item Une borne \terme{must} (sous-approximation) : \og $C \supseteq
    Y$\fg{}, ici \og change \emph{à chaque} rendu\fg{}. C'est un fait de
    certitude. Les règles s'en servent pour \emph{tirer}, jusqu'au niveau
    \Error{}.
\end{itemize}
Chaque borne prend trois valeurs, ordonnées : \code{Never} ($C = \emptyset$),
\code{OnSet} ($\mathit{Set}(S)$), \code{Every} (tous les rendus). Une paire
(may, must) cohérente vérifie must $\le$ may. Le \cref{fig:statevalue-stabilite-maymust}
place les six paires possibles.

\begin{figure}[htbp]\centering
\begin{tikzpicture}[x=34mm,y=17mm]
  % axes
  \draw[fleche] (-0.62,-0.5) -- (2.6,-0.5);
  \node[etiquette] at (1,-1.0) {borne may : $C \subseteq \dots$};
  \draw[fleche] (-0.62,-0.5) -- (-0.62,2.5);
  \node[etiquette] at (-0.62,2.75) {borne must : $C \supseteq \dots$};
  \foreach \x/\n in {0/\code{Never},1/\code{OnSet},2/\code{Every}} \node[etiquette,below] at (\x,-0.52) {\n};
  \foreach \y/\n in {0/\code{Never},1/\code{OnSet},2/\code{Every}} \node[etiquette,left] at (-0.64,\y) {\n};
  % points retenus
  \node[bloc,align=center] (s) at (0,0) {\code{Stable}\\{\scriptsize \code{useRef}, setter, littéral}};
  \node[bloc,align=center] (v) at (1,0) {\code{Versioned(S)}\\{\scriptsize lecture d'un état objet}};
  \node[bloc,align=center] (u) at (2,0) {\code{Unknown}\\{\scriptsize \code{c ? {} : obj}}};
  \node[bloc,align=center] (p) at (2,2) {\code{PerRender}\\{\scriptsize \code{{}} dans le rendu}};
  % points écartés
  \node[cfgbloc,dashed,align=center,text=ra-gray] (oo) at (1,1) {(OnSet, OnSet)\\{\scriptsize change exactement aux sets}};
  \node[cfgbloc,dashed,align=center,text=ra-gray] (eo) at (2,1) {(Every, OnSet)\\{\scriptsize aux sets, peut-être plus}};
  \node[etiquette,align=center] at (0.35,1.75) {must $>$ may :\\incohérent};
\end{tikzpicture}
\caption{Le carré may/must de l'ADR-017. Chaque point est une paire (borne
may, borne must) sur la trace de changement $C$ ; seules les paires avec
must $\le$ may ont un sens. Les quatre points encadrés en bleu sont retenus
(avec $\Bot$, non figuré) ; les deux points en pointillés, \og doit changer
aux sets\fg{}, sont écartés. Le join affaiblit les deux bornes : il va vers
la droite et vers le bas.}
\label{fig:statevalue-stabilite-maymust}
\end{figure}

Les deux points écartés n'avaient qu'un consommateur prospectif, la
détection du churn. Or celle-ci a de toute façon besoin d'une analyse
d'atteignabilité du corps de l'effet (le set a-t-il lieu sur tous les
chemins ?), qui n'est pas une propriété de valeur. Elle vit donc dans la
règle (\cref{chap:infinite-loop}) et le graphe de churn
(\cref{chap:churn}), pas dans le treillis. L'ADR-017 garde le produit
complet comme cadre d'extension si un second consommateur apparaît.

\subsection{Le treillis \code{Stability}}

\begin{definition}{Treillis de stabilité versionnée}{stabilite}
  Le treillis \code{Stability} a pour éléments $\Bot$, \code{Stable},
  $\code{Versioned}(S)$ pour tout ensemble $S$ de slots qualifiés avec
  $1 \le |S| \le 4$, \code{VersionedTop}, \code{PerRender} et
  \code{Unknown} ($\Top$). Sa concrétisation, en termes de traces de
  changement (\cref{def:statevalue-stabilite-trace}) :
  \[
  \begin{array}{rcl}
    \gam(\Bot) &=& \emptyset\\
    \gam(\code{Stable}) &=& \{\, \text{traces avec } C = \emptyset \,\}\\
    \gam(\code{Versioned}(S)) &=& \{\, \text{traces avec } C \subseteq \mathit{Set}(S) \,\}\\
    \gam(\code{VersionedTop}) &=& \{\, \text{traces avec } C \subseteq \mathit{Set}(S) \text{ pour un certain ensemble } S \,\}\\
    \gam(\code{PerRender}) &=& \{\, \text{traces avec } C = \{1, \dots, n\} \,\}\\
    \gam(\code{Unknown}) &=& \text{toutes les traces}
  \end{array}
  \]
  L'ordre est $\Bot \lleq \code{Stable} \lleq \code{Versioned}(S) \lleq
  \code{Versioned}(T) \lleq \code{VersionedTop} \lleq \code{Unknown}$ pour
  $S \subseteq T$, et $\Bot \lleq \code{PerRender} \lleq \code{Unknown}$ ;
  \code{PerRender} est incomparable avec \code{Stable}, les
  \code{Versioned} et \code{VersionedTop}.
\end{definition}

Les quatre noms se lisent sur le carré : \code{Stable} est (Never, Never),
\code{Versioned} est (OnSet, Never), \code{PerRender} est (Every, Every),
\code{Unknown} est (Every, Never). \code{Versioned} est une borne may ;
\code{PerRender} est une borne must. Le code transcrit la définition, avec un
commentaire qui reproduit le diagramme de Hasse :

\rustsnippet[label={lst:statevalue-stabilite-enum}]{src/domains/impls/stability.rs}{9}{52}{le treillis de stabilité versionnée}

Un slot qualifié \code{QualifiedSlot} est un couple \code{(ComponentId,
HookLabel)} (\src{src/ir/types.rs}{6}{9}) : le label d'un hook n'est unique
qu'au sein de son composant, et un \code{Versioned} peut traverser les
frontières de composants par les props. L'ADR-017 écrit encore
\code{(Symbol, HookLabel)} ; le passage à \code{ComponentId} date de
\issue{7}. Le \cref{fig:statevalue-stabilite-hasse} redessine le diagramme.

\begin{figure}[htbp]\centering
\begin{tikzpicture}[treillis,x=15mm,y=12mm]
  \node (top) at (1.5,5) {\code{Unknown} ($\Top$)};
  \node (vtop) at (0,4) {\code{VersionedTop}};
  \node (v4) at (0,3.1) {\code{Versioned({a,b,c,d})}};
  \node (vdots) at (0,2.45) {$\vdots$};
  \node (vab) at (-1,1.8) {\code{Versioned({a,b})}};
  \node (vac) at (1,1.8) {\code{Versioned({a,c})}};
  \node (va) at (-1,0.9) {\code{Versioned({a})}};
  \node (vb) at (0.2,0.9) {\code{{b}}};
  \node (vc) at (1.2,0.9) {\code{{c}}};
  \node (st) at (0,0) {\code{Stable}};
  \node (per) at (3.3,2.2) {\code{PerRender}};
  \node (bot) at (1.5,-1) {$\Bot$};
  \draw (top) -- (vtop); \draw (top) -- (per);
  \draw (vtop) -- (v4); \draw (v4) -- (vdots);
  \draw (vdots) -- (vab); \draw (vdots) -- (vac);
  \draw (vab) -- (va); \draw (vab) -- (vb.north); \draw (vac) -- (va); \draw (vac) -- (vc);
  \draw (va) -- (st); \draw (vb) -- (st); \draw (vc) -- (st);
  \draw (st) -- (bot); \draw (per) -- (bot);
  \node[etiquette,align=left] at (-2.6,2.3) {borne may :\\ chaîne $\subseteq$,\\ $|S| \le 4$};
  \node[etiquette,align=left] at (4.3,2.2) {borne\\ must};
\end{tikzpicture}
\caption{Diagramme de Hasse de \code{Stability} (d'après le commentaire de
\src{src/domains/impls/stability.rs}{20}{30}). À gauche, la chaîne may
ordonnée par inclusion des ensembles de slots, bornée par
\code{VERSIONED_LABELS_THRESHOLD = 4} ; à droite, le fait must
\code{PerRender}. Les deux branches ne se rejoignent qu'en $\Top$ et en
$\Bot$. Les étiquettes \code{{b}} et \code{{c}} abrègent
\code{Versioned({b})} et \code{Versioned({c})}.}
\label{fig:statevalue-stabilite-hasse}
\end{figure}

\subsection{Canonisation}

\rustsnippet[label={lst:statevalue-stabilite-versioned}]{src/domains/impls/stability.rs}{54}{70}{les constructeurs canonisants}

\begin{invariant}{Forme canonique de \code{Versioned}}{statevalue-stabilite-versioned}
  Un \code{Versioned(S)} a toujours $1 \le |S| \le 4$. \og Versionné par
  rien\fg{} veut dire \og ne change jamais\fg{} : $\code{Versioned}(\emptyset)
  \equiv \code{Stable}$ ; au-delà de quatre slots, la valeur est élargie à
  \code{VersionedTop}. L'invariant est garanti par \code{versioned} (qu'emploient
  \code{join} et \code{meet}) et par \code{versioned_by} (un singleton). Le
  variant étant public, rien n'empêche un appelant de construire
  \code{Versioned} à la main ; aucune assertion ne le vérifie.
\end{invariant}

Le seuil reprend le motif de \code{StrConst} (\cref{sec:treillis-simples-strconst}) :
un ensemble borné, élargi à son $\Top$ propre au-delà de quatre éléments.

\subsection{Ordre, join, meet, widen}

\rustsnippet[label={lst:statevalue-stabilite-ord-stab}]{src/domains/impls/stability.rs}{72}{101}{l'ordre de \code{Stability}}

L'ordre se lit de haut en bas : l'égalité d'abord, puis $\Bot$ et
\code{Unknown} aux extrémités, puis la chaîne may (\code{Stable} sous tout
\code{Versioned}, inclusion stricte des ensembles, tout \code{Versioned} sous
\code{VersionedTop}) ; tout le reste, c'est-à-dire \code{PerRender} face à
la chaîne may, et deux ensembles non inclus l'un dans l'autre, est
incomparable (\code{None}).

\rustsnippet[label={lst:statevalue-stabilite-join-stab}]{src/domains/impls/stability.rs}{108}{146}{join, meet et widen de \code{Stability}}

Ces méthodes sont \emph{inhérentes} ; l'implémentation du trait
\code{AbstractDomain} leur délègue
(\src{src/domains/impls/stability.rs}{149}{168}). Comme Rust préfère la méthode inhérente à la méthode de trait
à la résolution de \code{self.join(other)}, il n'y a pas de récursion.

\begin{theoreme}{Soundness du join de \code{Stability}}{statevalue-stabilite-join-stab}
  Pour tous $a, b$ : $\gam(a) \cup \gam(b) \subseteq \gam(a \join b)$.
\end{theoreme}
\begin{proof}
  Par cas, dans l'ordre des bras du code.
  \emph{Égalité} : $a \join a = a$. \emph{$\Bot$} : $\gam(\Bot) =
  \emptyset$, donc $\gam(\Bot) \cup \gam(x) = \gam(x)$.
  \emph{\code{Unknown}} : $\gam(\code{Unknown})$ contient toutes les traces.
  \emph{\code{Stable} et \code{Versioned}(S)} : une trace avec $C =
  \emptyset$ vérifie $C \subseteq \mathit{Set}(S)$ ; même argument avec
  \code{VersionedTop}. \emph{\code{Versioned}(S) et \code{Versioned}(T)} :
  $\mathit{Set}$ est croissant, donc $C \subseteq \mathit{Set}(S)$ ou $C
  \subseteq \mathit{Set}(T)$ entraîne $C \subseteq \mathit{Set}(S \cup T)$ ;
  si $|S \cup T| > 4$, \code{versioned} rend \code{VersionedTop}, dont la
  concrétisation contient celle de $\code{Versioned}(S \cup T)$ ; et
  $S \cup T$ n'est jamais vide. \emph{\code{VersionedTop} et
  \code{Versioned}} : immédiat. \emph{Dernier bras} (\code{PerRender} face à
  la chaîne may) : le résultat est \code{Unknown}, qui contient tout.
\end{proof}

\begin{proposition}{Borne supérieure, hauteur et terminaison}{statevalue-stabilite-lub}
  Sous l'\cref{inv:statevalue-stabilite-versioned}, $a \join b$ est la plus
  petite borne supérieure de $a$ et $b$ dans l'ordre de
  \cref{lst:statevalue-stabilite-ord-stab}. Toute chaîne strictement
  croissante a au plus 8 éléments ($\Bot \sqsubset \code{Stable} \sqsubset$
  quatre \code{Versioned} de tailles 1 à 4 $\sqsubset \code{VersionedTop}
  \sqsubset \code{Unknown}$), donc \code{widen = join} termine.
\end{proposition}
\begin{proof}
  Seuls deux cas ne sont pas évidents. $\code{Stable} \join
  \code{PerRender}$ : le seul majorant commun est \code{Unknown}, puisque
  \code{PerRender} n'est sous rien d'autre. $\code{Versioned}(S) \join
  \code{Versioned}(T)$ avec $|S \cup T| > 4$ : un majorant
  $\code{Versioned}(U)$ exigerait $U \supseteq S \cup T$, donc $|U| > 4$, ce
  que l'invariant interdit ; restent \code{VersionedTop} et \code{Unknown},
  et le plus petit est \code{VersionedTop}.
\end{proof}

\begin{piege}
  $\code{Versioned}(S) \meet \code{PerRender} = \Bot$ est bien la borne
  inférieure \emph{dans l'ordre abstrait}, mais les concrétisations
  s'intersectent : une trace où chaque rendu est précédé d'un set d'un slot
  de $S$ change à chaque rendu \emph{et} seulement aux sets. Le meet est
  donc strictement plus petit que l'intersection des concrétisations ;
  l'utiliser pour raffiner une valeur ferait disparaître des comportements
  réels, un FN. Il n'est appelé par aucun code de production au commit
  \snapshotCommit{} ; ne pas s'en servir sans revoir ce point.
\end{piege}

\begin{decision}[ADR-017 --- stabilité versionnée, bornes may et must]
\textbf{Contexte} : les cinq FP F5 sur des états objets, et le FN
\code{ObjChurn} que leur correction naïve aurait ouvert.
\textbf{Décide} : (1) le treillis ci-dessus ; (2) une conversion \og côté
lecture\fg{} en un seul point, la lecture d'un slot
(\cref{sec:statevalue-stabilite-lecture}) ; (3) un bras churn dans
\regle{infinite-loop} (\Error{} si la dep \emph{est} le slot, que le set a
lieu sur tous les chemins et que la valeur écrite est \code{PerRender} ;
\Warning{} si la dep est \code{Versioned(S)} avec le slot dans $S$) ; (4) le
recâblage des consommateurs (\regle{always-unstable-deps} ne tire plus que
sur \code{PerRender}). \textbf{Alternatives refusées} : encoder un compteur
d'événements dans le store pour faire tirer le widening (\og non-standard
hack\fg{}) ; un nouveau nom de diagnostic (même classe de panne) ; garder les
deux points \og doit changer aux sets\fg{} (un seul consommateur, qui a de
toute façon besoin d'atteignabilité). \textbf{Coût accepté} :
\code{Stability} n'est plus \code{Copy} (il possède un \code{BTreeSet}) ; les
erreurs de compilation ont forcé l'audit de tous les consommateurs.
\end{decision}

% ===========================================================================
\section{Deux vues d'un même état : la conversion côté lecture}
\label{sec:statevalue-stabilite-lecture}

\subsection{Le problème}

Le treillis sait maintenant dire \og ne change qu'aux sets de
\code{ctx}\fg{}. Reste à savoir \emph{où} ce fait apparaît. Dans
\code{ObjectStateDep}, le store du slot reçoit deux écritures : le
littéral d'initialisation et celui du handler, deux objets neufs. Son contenu
convergé est donc, à juste titre, \code{ref(PerRender)} : chaque valeur
\emph{écrite} est une allocation fraîche. La sonde le confirme :
\begin{sortie}
-- ObjectStateDep
   state[0] = ref(PerRender)   widened=false  num=Interval { lo: inf, hi: -inf, is_int: true }
\end{sortie}
Si la lecture \code{ctx} rendait cette valeur telle quelle, on retomberait
sur le FP F5. Il faut que ce qu'un rendu \emph{lit} dise autre chose que ce
que les événements \emph{écrivent}.

\subsection{Le code : un seul point de conversion}

Le fait React qui tranche est celui-ci : un slot d'état ne change \emph{que}
par son setter. Toute lecture de \code{StateVal(l)} a donc la borne may
\og ne change qu'aux sets de $l$\fg{}, quelles que soient les valeurs écrites.
L'évaluateur le traduit à un seul endroit :

\rustsnippet[label={lst:statevalue-stabilite-stateval}]{src/domains/transfer/state_value.rs}{122}{135}{la conversion côté lecture (ADR-017)}

Si la composante \code{reference} de la valeur stockée est peuplée, elle est
remplacée par \code{Versioned({(composant, l)})} ; les autres composantes
passent inchangées. Le champ \code{ctx.component} de l'\code{AnalysisCtx}
(\cref{sec:statevalue-stabilite-contextes}) fournit le propriétaire du slot.
Une lecture dont la composante \code{reference} est $\Bot$ (un état
purement numérique, par exemple) n'est pas convertie : elle ne porte aucune
étiquette de version, ce qui aura une conséquence au
\cref{sec:statevalue-stabilite-limites}. Tout ce qui est en aval hérite de
la conversion gratuitement : liaisons d'environnement
(\code{const o = obj}), évaluation des deps, props passées à un enfant.

Le \cref{chap:react-modele} a énoncé l'invariant qui en découle, la
\emph{double vue de l'état} (\cref{inv:react-modele-double-vue}) ; le
\cref{fig:statevalue-stabilite-double-vue} le met en image.

\begin{figure}[htbp]\centering
\begin{tikzpicture}[node distance=7mm and 9mm]
  \node[bloc,align=center] (w1) {\code{useState({ locale: "en" })}\\{\scriptsize amorçage}};
  \node[bloc,align=center,below=of w1] (w2) {\code{setCtx({ locale: "fr" })}\\{\scriptsize handler}};
  \coordinate (mid) at ($(w1.east)!0.5!(w2.east)$);
  \node[bloc,align=center,right=12mm of mid,minimum height=16mm] (store) {\textbf{store} (vue événement)\\ join des valeurs écrites\\ \code{ref(PerRender)}};
  \node[bloc,align=center,right=18mm of store,minimum height=16mm] (read) {\textbf{lecture} \code{StateVal(l)}\\ (vue inter-rendus)\\ \code{ref(Versioned({(C,l)}))}};
  \draw[fleche] (w1) -- (store); \draw[fleche] (w2) -- (store);
  \draw[flechemust] (store) -- node[etiquette,above]{conversion} node[etiquette,below]{\code{versioned_by}} (read);
  \node[cfgbloc,align=left,below=12mm of store] (r1) {\regle{redundant-set-state}\\ \code{value_freshness} (churn)};
  \node[cfgbloc,align=left,below=12mm of read] (r2) {deps : \regle{always-unstable-deps},\\ \regle{infinite-loop}, \regle{missing-deps},\\ mémos, props};
  \draw[flechemay] (store) -- node[etiquette,left]{lit} (r1);
  \draw[flechemay] (read) -- node[etiquette,right]{lit} (r2);
\end{tikzpicture}
\caption{La double vue de l'état sur \code{ObjectStateDep}. Le store range
le join des valeurs \emph{écrites}, deux allocations fraîches ; toute
lecture du slot rend la vue \emph{inter-rendus}, ce que \code{Object.is}
compare d'un rendu à l'autre. Les consommateurs qui comparent des valeurs
écrites lisent le store ; ceux qui raisonnent sur les deps lisent
l'évaluation.}
\label{fig:statevalue-stabilite-double-vue}
\end{figure}

\begin{proposition}{Soundness de la conversion côté lecture}{statevalue-stabilite-lecture}
  Sous l'hypothèse (H1) \og aucun setter n'est appelé pendant le rendu\fg{},
  la trace de changement de toute lecture du slot $l$ vérifie $C \subseteq
  \mathit{Set}(\{l\})$ : la composante $\code{Versioned}(\{l\})$ est une
  borne may correcte.
\end{proposition}
\begin{proof}
  La valeur d'un état au rendu $r_i$ est la valeur de son slot, que seuls
  les appels à son setter modifient (en file d'attente, traités au rendu
  suivant). Si $r_i$ n'est précédé, depuis $r_{i-1}$, d'aucun appel au setter
  de $l$, et qu'aucun appel n'a eu lieu pendant $r_{i-1}$ lui-même (H1), la
  valeur est la même référence, et \code{Object.is} réussit : $i \notin C$.
\end{proof}

L'ADR-017 nomme les deux hypothèses dont dépend cette conversion, et elles
sont couvertes chacune à part :
\begin{enumerate}
  \item \emph{H1, les sets ont lieu hors rendu.} Un setter appelé pendant le
    rendu avec une référence fraîche fait vraiment changer le slot à chaque
    rendu, ce que \code{Versioned} sous-estime. Cette violation a son propre
    diagnostic, \regle{setter-in-render}, de niveau \Error{} : c'est la
    soundness en couches du \cref{chap:react-modele}
    (\cref{def:soundness-couches}).
  \item \emph{H2, le couplage avec le bras churn.} \regle{infinite-loop}
    accepte une dep \code{Versioned} comme \og garde\fg{} qui empêche
    l'effet de tourner à chaque rendu. C'est sound \emph{seulement} parce
    que le bras churn traite le cas où l'effet écrit lui-même une référence
    fraîche dans le slot qui le versionne. Retirer l'un sans l'autre rouvre
    le FN \code{ObjChurn}.
\end{enumerate}

\begin{exemple}{Les cinq composants de \file{stability.tsx}}{statevalue-stabilite-cinq}
Le fichier \file{/tmp/ch10/stability.tsx} réunit \code{ObjectStateDep},
\code{ObjChurn} et trois composants de la \cref{sec:statevalue-stabilite-memo}
(\code{InlineObjectDep}, \code{MemoizedDep}, \code{CallbackDep}) :
\begin{sortie}
$ reactant check /tmp/ch10/stability.tsx --show-clean
  CallbackDep  (4 hooks)  /tmp/ch10/stability.tsx  ✓
  InlineObjectDep  (1 hooks)  /tmp/ch10/stability.tsx
    warn   always-unstable-deps  [hook:0]  (line 5:2)  this effect depends on `opts`, a new reference every render, so `Object.is` always differs and the effect re-runs on every render regardless of the other deps
       (1 trace step(s), rerun with --trace)
  MemoizedDep  (2 hooks)  /tmp/ch10/stability.tsx  ✓
  ObjChurn  (2 hooks)  /tmp/ch10/stability.tsx
    error  infinite-loop  [hook:0]  (line 30:2)  this effect recreates object state `obj` it depends on. Every run stores a fresh reference (`Object.is` always fails) and re-triggers itself: infinite render loop
       (1 trace step(s), rerun with --trace)
  ObjectStateDep  (3 hooks)  /tmp/ch10/stability.tsx  ✓

⚠  1 error(s), 1 warning(s) across 1 file(s).
\end{sortie}
\code{ObjectStateDep} est propre : sa dep lit \code{Versioned}, et
\regle{always-unstable-deps} exige \code{PerRender}. Le FP F5 est mort, et
\code{ObjChurn} est passé d'un \Warning{} noyé dans les FP à une \Error{}.
\end{exemple}

\subsection{Le churn d'objet, certifié par la structure}

Sur \code{ObjChurn}, le store vaut aussi \code{ref(PerRender)}, et la lecture
\code{obj} vaut \code{Versioned({(C,0)})}. Rien ne s'élargit (sonde :
\code{widened=false}), donc le bras point fixe d'\regle{infinite-loop} ne peut
rien voir. La certitude ne vient pas non plus de \code{Versioned} : c'est une
borne may (\og ne change \emph{qu'aux} sets\fg{}), pas \og change à chaque
set\fg{}. Elle vient de la \emph{structure} :
\begin{itemize}
  \item la dep \code{obj} \emph{est} le slot, sans accès de champ
    (must-change dès que la valeur écrite est fraîche) ;
  \item la valeur écrite \code{{...obj, b: 2}} est un \code{ObjectLit}, donc
    \code{PerRender} au sens must, et l'écriture a lieu sur tous les chemins du
    corps de l'effet (must-reach) ;
  \item React ré-exécute un effet dont une dep a changé (must-rerun).
\end{itemize}
Trois faits must, aucun fait may : une \Error{} (\cref{chap:infinite-loop}).
Côté moteur, c'est \code{value_freshness} qui classe la valeur écrite pour
le graphe de churn (\cref{chap:churn}) ; elle lit la vue événement :

\rustsnippet[label={lst:statevalue-stabilite-freshness}]{src/engine/written.rs}{173}{200}{la fraîcheur d'une valeur écrite, lue sur la composante référence}

Le commentaire d'en-tête dit l'essentiel : le churn ne concerne que la
sorte référence, parce qu'une valeur numérique élargie (\code{count + 1})
change mais n'échoue jamais \code{Object.is} \og par fraîcheur\fg{}. Une
valeur versionnée par le slot cible seul est son propre contenu : la
réécrire n'est pas un changement (\issue{155}).

\begin{piege}
  L'inverse n'est pas vrai : une valeur versionnée par le slot cible mais
  \emph{dérivée} de lui (un mémo \code{useMemo(() => ({ n: s.n + 1 }), [s])} réécrit dans \code{s}) est lue \og non fraîche\fg{}, alors qu'elle
  change exactement quand le cycle tourne. C'est un FN connu, ouvert sous
  l'étiquette \code{soundness-bug} (\issue{157}) ; sa correction demande un
  point \og exactement le contenu du slot $q$\fg{} \emph{sous}
  $\code{Versioned}(\{q\})$ dans le treillis, dégradé en \code{Versioned}
  par toute transformation. En attendant, la lecture d'ADR-017 tient.
\end{piege}

\subsection{Lire un champ d'un objet versionné}

Un accès \code{obj.field} commence par le tas : si le récepteur désigne des
objets alloués connus, la valeur est le join des champs lus
(\cref{chap:transfert-stores}). Sinon, l'évaluateur garde les étiquettes de
version du récepteur plutôt que de tomber à $\Top$ :

\rustsnippet[label={lst:statevalue-stabilite-field}]{src/domains/transfer/state_value.rs}{617}{628}{un champ d'objet versionné garde les étiquettes}

La \emph{sorte} du champ reste inconnue (toutes les autres composantes sont à
$\Top$), mais sa borne may survit : un champ d'un état objet ne peut changer
qu'aux sets de cet état, sous la même hypothèse \og pas de mutation pendant
le rendu\fg{} que la lecture d'un slot (les mutations en place ont leur
diagnostic, \regle{state-mutation}). Une dep \code{obj.field} ne peut
toutefois donner qu'un \Warning{} dans le bras churn :
\code{setObj({...obj, other: 2})} préserve la référence du champ, donc la
dep peut ne pas changer.

% ===========================================================================
\section{Littéraux, \code{useRef}, \code{useMemo}, \code{useCallback}}
\label{sec:statevalue-stabilite-memo}

\subsection{Où naît une référence}

La composante \code{reference} n'est peuplée que par un petit nombre de
formes, toutes dans l'évaluateur :

\rustsnippet{src/domains/transfer/state_value.rs}{153}{156}{les littéraux alloués et les éléments natifs}

\rustsnippet{src/domains/transfer/state_value.rs}{176}{181}{appels et \code{new}}

Le \cref{tab:statevalue-stabilite-sources} rassemble ces sources.

\begin{table}[htbp]\centering\small
\begin{tabular}{p{47mm}p{33mm}p{62mm}}
\hline
\textbf{Expression} & \textbf{Valeur} & \textbf{Raison} \\ \hline
\code{{...}}, \code{[...]}, \code{() => ...} & \code{ref(PerRender)} & allocation à chaque évaluation (fait must) \\
\code{new X(...)} & \code{ref(PerRender)} & idem, depuis \code{e67b10a} (\issue{158}) \\
\code{arr.map(f)}, \code{filter}, \code{flat}, \code{split}, \code{structuredClone(x)}, \code{JSON.parse}\dots & \code{ref(PerRender)} & \code{returns_fresh_reference} (\src{src/domains/transfer/state_value.rs}{238}{259}) \\
\code{s.slice(...)}, \code{concat} & $\Top$ & peut rendre une chaîne, donc une primitive (\issue{22}) \\
autre appel & $\Top$ & rien n'est connu du résultat \\
élément JSX natif & \code{ref(Stable)} & non documenté, un FN
(\cref{sec:statevalue-stabilite-nativeelem}) \\
\code{useRef()} & \code{ref(Stable)} & même conteneur à chaque rendu \\
\code{setX} (setter) & \code{setter(...)} & identité garantie par React \\
lecture \code{StateVal(l)} & \code{Versioned({l})} sur \code{reference} & conversion côté lecture \\
\code{useMemo} / \code{useCallback} & \code{ref(...)} & \code{recompute_memo}, ci-dessous \\
\hline
\end{tabular}
\caption{Les sources de la composante \code{reference}.}
\label{tab:statevalue-stabilite-sources}
\end{table}

\code{useRef} passe par un marqueur de hook dont le commentaire explique
pourquoi il est une référence et pas \code{undefined} :

\rustsnippet{src/domains/transfer/state_value.rs}{142}{147}{les marqueurs de hooks}

Le commentaire qui précède (\src{src/domains/transfer/state_value.rs}{137}{141})
mérite d'être retenu : lire le résultat
d'un hook personnalisé non résolu comme \code{undefined} le rendait
\emph{prouvablement stable} (\code{to_stability} joint \code{Stable} pour
\code{undef}) et faisait taire toutes les règles de stabilité, un FN. Il vaut
donc $\Top$.

\subsection{La valeur d'un mémo : \code{recompute_memo}}

\code{useMemo(f, deps)} et \code{useCallback(f, deps)} rendent la même
référence tant que leurs deps ne changent pas. Leur stabilité est donc le
join des stabilités de leurs deps. Le moteur la recalcule par la méthode
\code{recompute_memo} du trait \code{Transfer} :

\rustsnippet[label={lst:statevalue-stabilite-memo}]{src/domains/transfer/state_value.rs}{56}{101}{la valeur d'un mémo, stabilité de ses deps}

Étape par étape :
\begin{enumerate}
  \item pas de tableau lisible (\code{deps.list()} vaut \code{None}) :
    \code{Unknown}. \code{Stable} serait une affirmation must sur une liste
    que l'IR n'a jamais vue ;
  \item un tableau \emph{connu} vide (\code{Arity::Exact(0)},
    \cref{sec:ir-deps}) : \code{Stable}, le mémo est calculé une fois ;
  \item un tableau dont toutes les entrées sont des spreads illisibles :
    \code{Unknown} ;
  \item sinon, un pli (fold) par join sur les deps, depuis $\Bot$ : une dep
    qui \emph{est} syntaxiquement \code{StateVal(l)} donne
    \code{versioned_by(component, l)} quelle que soit la sorte du slot ;
    toute autre dep est évaluée normalement, et contribue sa composante
    versionnée si elle en a une, sa projection \code{to_stability} sinon.
\end{enumerate}
Le résultat est rangé dans la composante \code{reference} d'une valeur dont
tout le reste est $\Bot$.

\begin{exemple}{Trois mémos}{statevalue-stabilite-memos}
  Sur \file{stability.tsx}, \code{MemoizedDep} écrit
  \code{useMemo(() => ({ a: 1 }), [])} ; \code{CallbackDep} écrit
  \code{useCallback(() => console.log(q), [q])} avec \code{q} un état objet
  mis à jour par un handler ; et sur \file{memo.tsx}, \code{MemoOverObjState}
  écrit \code{useMemo(() => ({ v: o.a }), [o])}. La sonde :
\begin{sortie}
-- MemoizedDep
   memo[0] (Memo) = ref(Stable)
-- CallbackDep
   state[0] = ref(PerRender)   widened=false  num=Interval { lo: inf, hi: -inf, is_int: true }
   memo[1] (Callback) = ref(Versioned({(ComponentId(4294967295), 0)}))
-- MemoOverObjState
   state[0] = ref(PerRender)   widened=false  num=Interval { lo: inf, hi: -inf, is_int: true }
   memo[1] (Memo) = ref(Versioned({(ComponentId(4294967295), 0)}))
\end{sortie}
  Les deux derniers illustrent la double vue : le store de l'état objet
  vaut \code{ref(PerRender)}, mais la dep lue est versionnée, et le mémo
  hérite de l'étiquette. Un effet qui dépend de \code{onLoad} ou de \code{a}
  est donc \og ne change qu'aux sets du slot 0\fg{}, et les trois composants
  sont propres.
\end{exemple}

La propagation des étiquettes le long d'une chaîne de mémos est gratuite :
$\code{Versioned}(\{X\}) \join \code{Stable} = \code{Versioned}(\{X\})$.
L'ADR-017 exige seulement que \code{to_stability} reste hors de ce chemin,
parce qu'il effacerait les étiquettes dès qu'une autre composante est
$\Top$ : c'est le rôle de \code{versioned_reference}
(\cref{sec:statevalue-stabilite-projections}).

\begin{piege}
  La valeur d'un mémo n'est \emph{que} sa stabilité, logée dans la
  composante \code{reference}. La sorte réelle du résultat (un nombre, un
  objet) est perdue : \code{useMemo(() => 1, [])} vaut \code{ref(Stable)} et
  non le nombre $1$. Avec une dep fraîche, la perte devient visible :
\begin{lstlisting}[language=TypeScript,caption={Un mémo primitif sur une dep fraîche (\file{/tmp/ch10/memo.tsx}, extrait)},label={lst:statevalue-stabilite-memoprim}]
export function MemoPrimitiveOverFreshDep() {
  const cfg = { k: 1 };
  const one = useMemo(() => 1, [cfg]);
  useEffect(() => { console.log(one); }, [one]);
  return <div>{one}</div>;
}
\end{lstlisting}
\begin{sortie}
  MemoPrimitiveOverFreshDep  (2 hooks)  /tmp/ch10/memo.tsx
    warn   always-unstable-deps  [hook:0]  (line 19:8)  this memo depends on `cfg`, a new reference every render, so `Object.is` always differs and the memo re-runs on every render regardless of the other deps
       → the value flows through `cfg`, bound here (line 18:8)
    warn   always-unstable-deps  [hook:1]  (line 20:2)  this effect depends on `one`, a new reference every render, so `Object.is` always differs and the effect re-runs on every render regardless of the other deps
       → the value flows through `one`, bound here (line 19:8)
\end{sortie}
  Le premier \Warning{} est juste : le mémo est recalculé à chaque rendu.
  Le second est un FP : \code{one} vaut toujours \code{1}, comparé par
  valeur, et l'effet ne tourne qu'une fois. La sonde donne
  \code{memo[0] (Memo) = ref(PerRender)}.
\end{piege}

% ===========================================================================
\section{Projections vers les règles}
\label{sec:statevalue-stabilite-projections}

\subsection{Quatre questions différentes}

Les règles ne lisent pas une \code{StateValue} en entier ; elles lui posent
des questions fermées, et chaque question a sa polarité :
\begin{itemize}
  \item \og cette dep est-elle, à coup sûr, une référence neuve à chaque
    rendu ?\fg{} (must ; \regle{always-unstable-deps}) ;
  \item \og cette valeur est-elle prouvablement stable, donc omissible des
    deps ?\fg{} (une preuve de sûreté ; \regle{missing-deps}) ;
  \item \og cette écriture peut-elle croître sans borne ?\fg{} (may ; le bras
    point fixe d'\regle{infinite-loop}) ;
  \item \og cet état ne prend-il que quelques valeurs primitives ?\fg{}
    (\regle{wasted-subtree-render}).
\end{itemize}
Le code fournit une projection par question. La plus ancienne, et la plus
délicate, est \code{to_stability}.

\subsection{\code{to_stability} : la règle \og motion-wins\fg{}}

\rustsnippet[label={lst:statevalue-stabilite-to-stab}]{src/domains/impls/state_value.rs}{320}{369}{la projection vers la stabilité}

L'algorithme :
\begin{enumerate}
  \item le résidu \code{other} force \code{Unknown} : une valeur opaque ne
    revendique jamais une (in)stabilité définie ;
  \item chaque composante peuplée contribue : un point numérique, un booléen
    exact, une chaîne singleton, \code{null}, \code{undefined}, un setter
    contribuent \code{Stable} ; un booléen $\Top$ ou plusieurs chaînes
    contribuent \code{Unknown} ; la référence contribue sa propre stabilité,
    sauf \code{PerRender}, qui lève le drapeau \og en mouvement\fg{} ;
  \item deux sortes peuplées ou plus donnent \code{Unknown} : deux sortes
    sont deux valeurs \code{Object.is}-distinctes, et la valeur peut passer
    de l'une à l'autre ;
  \item \terme{motion-wins} : un intervalle non ponctuel ou une référence
    \code{PerRender} rend \code{PerRender}, même si une autre composante a
    dit \code{Unknown}.
\end{enumerate}

La sonde, sur des valeurs construites à la main :
\begin{sortie}
null ⊔ 1 = number[1, 1]|null; to_stability = Unknown
widen -> number[1, inf]|null; to_stability = PerRender; is_unbounded=true
boolean ⊔ [0,5] to_stability = PerRender (motion wins over Unknown)
setter: setter(#4294967295#1) to_stability=Stable typeof=Some("function")
Versioned ⊔ Stable = Versioned({(ComponentId(4294967295), 0)})
Versioned ⊔ PerRender = Unknown
top-with-versioned-ref to_stability = Unknown; versioned_reference = Some(Versioned({(ComponentId(4294967295), 0)}))
\end{sortie}
La première ligne est la règle des sortes (deux points, deux sortes :
\code{Unknown}) ; la deuxième et la troisième sont motion-wins ; la dernière
montre pourquoi \code{versioned_reference} existe : sur une valeur dont
toutes les composantes sont $\Top$ sauf une référence versionnée,
\code{to_stability} rend \code{Unknown} et efface l'étiquette.

\begin{piege}
  \code{PerRender} n'a pas le même sens selon sa provenance. Dans la
  composante \code{reference}, c'est une borne \emph{must} : une nouvelle
  référence à chaque rendu, garantie. Rendu par \code{to_stability} pour un
  intervalle non ponctuel, c'est seulement \og \emph{peut} changer à chaque
  rendu\fg{} : l'ADR-017 l'assume (\og the name reads oddly for a number but
  means `may change every render', kind-agnostic\fg{}). Un consommateur qui
  a besoin du fait must doit lire la composante, pas la projection. Le
  message de \regle{missing-deps} montre la confusion au grand jour :
\begin{lstlisting}[language=TypeScript,caption={Un nombre à deux valeurs (\file{/tmp/ch10/limits.tsx}, extrait)},label={lst:statevalue-stabilite-branchy}]
export function Branchy({ flag }: { flag: boolean }) {
  const x = flag ? 1 : 2;
  useEffect(() => { console.log(x); }, []);
  return <div>{x}</div>;
}
\end{lstlisting}
\begin{sortie}
    warn   missing-deps  [hook:0]  var:x  (line 13:2)  `x` is used in this effect but not in its deps array, and it is recreated on every render
\end{sortie}
  L'avertissement est fondé (\code{x} change si \code{flag} change), mais
  \og recreated on every render\fg{} est faux : \code{x} vaut $[1,2]$, que
  \code{to_stability} projette sur \code{PerRender}, et
  \code{describe_value} (\src{src/rules/helpers/mod.rs}{36}{52}) traduit
  \code{PerRender} par cette phrase.
\end{piege}

\subsection{Les prédicats}

Le fait must a son propre prédicat, déjà rencontré
(\cref{lst:statevalue-stabilite-as-setter}) :
\code{is_unstable_reference_only} est vrai si la composante \code{reference}
vaut \code{PerRender} \emph{et} qu'aucune autre sorte n'est peuplée. C'est
celui qu'appellent \regle{always-unstable-deps}
(\srcline{src/rules/impls/always_unstable_deps.rs}{164}), le classement de
fraîcheur du churn (\srcline{src/engine/written.rs}{188}) et
\code{returns_verdict_of} (\srcline{src/rules/api/query.rs}{214}). Les
prédicats de croissance et de finitude :

\rustsnippet[label={lst:statevalue-stabilite-unbounded}]{src/domains/impls/state_value.rs}{371}{401}{croissance non bornée et finitude}

\code{is_unbounded} est la garde n°2 du bras point fixe
d'\regle{infinite-loop} : une borne numérique infinie, une référence
\code{PerRender} ou le résidu. Le test \code{!self.num.is_bottom()} n'est
pas décoratif : l'intervalle $\Bot$ a des bornes sentinelles infinies
(\cref{lst:statevalue-stabilite-bottom}), et sans ce test toute valeur non
numérique serait \og non bornée\fg{} (test
\code{is_unbounded_requires_active_num_slot}). \code{is_finitely_valued}
sert à \regle{wasted-subtree-render}
(\srcline{src/rules/impls/wasted_subtree_render.rs}{167}) : un état qui ne
prend que quelques primitives ne provoque un rendu qu'à ses
\emph{transitions}, puisque React abandonne une écriture
\code{Object.is}-égale. C'est le seul prédicat qui lise le bit
\code{is_int} hors raffinement.

\rustsnippet[label={lst:statevalue-stabilite-versioned-ref}]{src/domains/impls/state_value.rs}{403}{417}{garder les étiquettes, tester la stabilité}

\code{is_stable} est \og \code{to_stability} vaut \code{Stable}\fg{} ; il
sert de preuve de sûreté à \regle{missing-deps}
(\srcline{src/rules/impls/missing_deps.rs}{88}), \regle{redundant-set-state}
et \regle{unnecessary-rerender}.

\begin{exemple}{Deux sortes, une seule sorte}{statevalue-stabilite-crosskind}
\begin{lstlisting}[language=TypeScript,caption={Deux états optionnels (\file{/tmp/ch10/kinds.tsx}, extrait)},label={lst:statevalue-stabilite-crosskind}]
export function CrossKind() {
  const [x, setX] = useState<string>();
  useEffect(() => { console.log(x); }, []);
  return <button onClick={() => setX("data")}>{x}</button>;
}

export function NeverWritten() {
  const [x] = useState<string>();
  useEffect(() => { console.log(x); }, []);
  return <div>{x}</div>;
}
\end{lstlisting}
  La sonde donne \code{string{"data"}|undefined} pour le premier état et
  \code{undefined} pour le second. Le premier mêle deux sortes, chacune un
  point : \code{to_stability} rend \code{Unknown}, \code{is_stable} est faux,
  et \regle{missing-deps} tire à raison (\code{x} passe de \code{undefined} à
  \code{"data"}) :
\begin{sortie}
  CrossKind  (3 hooks)  /tmp/ch10/kinds.tsx
    warn   missing-deps  [hook:1]  var:x  (line 5:2)  `x` is used in this effect but not in its deps array, and its value may change between renders
\end{sortie}
  Le second n'a qu'une sorte, un point : \code{Stable}, et \code{NeverWritten}
  est propre. Aucune conversion \code{Versioned} n'a lieu, puisque la
  composante \code{reference} est $\Bot$.
\end{exemple}

\subsection{Côté règles : \code{StabilityVerdict}}

La couche règles (\cref{chap:api-regles}) n'expose pas \code{Stability}
directement, mais une projection totale qui replie $\Bot$ et $\Top$ sur le
côté may :

\rustsnippet[label={lst:statevalue-stabilite-verdict}]{src/rules/api/query.rs}{137}{179}{le verdict de stabilité côté règles}

Trois points de lecture. \code{Bottom} ne vaut pas \code{Stable} : un
$\Bot$ n'est pas une preuve de stabilité, il part avec \code{Unknown}.
\code{VersionedTop} devient un \code{Versioned} d'ensemble vide
(\og change à un slot inconnu\fg{}). Et la documentation de \code{PerRender}
répète la mise en garde du piège ci-dessus. Deux fonctions libres
s'appuient dessus : \code{stability_verdict_of}
(\srcline{src/rules/api/query.rs}{185}) et \code{may_change_of}
(\srcline{src/rules/api/query.rs}{225}), l'unique sonde \og peut
changer\fg{} sûre face à $\Top$ (\adr{21}).

\begin{table}[htbp]\centering\small
\begin{tabular}{p{40mm}p{22mm}p{76mm}}
\hline
\textbf{Projection} & \textbf{Polarité} & \textbf{Consommateurs principaux} \\ \hline
\code{is_unstable_reference_only} & must & \regle{always-unstable-deps}, \code{value_freshness}, \code{returns_verdict_of}, \file{helpers/jsx.rs}, \file{engine/fixpoint.rs} \\
\code{is_stable} & sûreté & \regle{missing-deps}, \regle{redundant-set-state}, \regle{unnecessary-rerender} \\
\code{is_unbounded} & may & \regle{infinite-loop} (bras point fixe) \\
\code{is_finitely_valued} & sûreté & \regle{wasted-subtree-render} \\
\code{to_stability} & mixte & \code{StabilityVerdict}, \code{describe_value}, \code{recompute_memo} \\
\code{versioned_reference} & may & \code{recompute_memo}, \code{eval_field_access} \\
\code{as_setter} & exacte & transfert des setters, \file{engine/setters.rs} \\
\hline
\end{tabular}
\caption{Les projections de \code{StateValue} et leurs consommateurs
(relevé \code{grep} au commit \snapshotCommit, dossier 04 §4.7).}
\label{tab:statevalue-stabilite-projections}
\end{table}

\begin{decision}[ADR-020 --- ne pas supprimer \code{to_stability}]
Lors de la campagne de dette technique, la suppression de
\code{to_stability} au profit de prédicats de première classe a été
proposée puis refusée (\adr{20}, non-changement n°4) : c'est une projection
légitime, documentée et testée (motion-wins, union de sortes vers
\code{Unknown}, setter stable), et aucun couplage avec perte d'information
démontré n'y cause de FN ou de FP. La supprimer serait du remaniement ou un
risque de soundness, pas une correction. La même campagne a refusé
d'extraire un combinateur \code{BoundedPowerset<T,N>} commun à
\code{StrConst} et \code{Stability} (non-changement n°5) : \code{Stability}
est un treillis plus riche que l'ensemble borné.
\end{decision}

% ===========================================================================
\section{Les contextes : la passerelle vers le moteur}
\label{sec:statevalue-stabilite-contextes}

\subsection{Ce dont une fonction de transfert a besoin}

Les fonctions de transfert ne manipulent pas seulement des
\code{StateValue} isolées. Pour évaluer \code{StateVal(l)}, il leur faut le
store d'état \emph{et} l'identité du composant, sans laquelle l'étiquette
\code{Versioned({(C, l)})} ne peut pas être fabriquée. Pour \code{MemoVal},
le store des mémos ; pour \code{obj.field}, le tas ; pour un appel à travers
une variable liée par \code{useCallback}, le CFG du corps ; pour
\code{<Child />}, toute la machinerie inter-composants. Le fichier
\file{src/domains/context.rs} regroupe ces besoins en quelques types, seule
flèche de dépendance \code{domains} $\to$ \code{engine} du périmètre.

\subsection{\code{AnalysisCtx} : le paquet d'état mutable}

\rustsnippet[label={lst:statevalue-stabilite-analysisctx}]{src/domains/context.rs}{105}{125}{le paquet d'état passé à chaque fonction de transfert}

L'environnement local n'y figure pas : il est passé à part, parce que
\code{eval_expr} le lit (\code{&}) alors que \code{exec_stmt} le modifie
(\code{&mut}). Le champ \code{component} n'est pas une \code{Option}, et son
commentaire dit pourquoi : une analyse est toujours l'analyse d'\emph{un}
composant, donc la provenance d'un slot (\code{Versioned}, \code{SetterVal})
porte toujours le vrai propriétaire ; une IR construite à la main et
analysée seule porte \code{ComponentId::SYNTHETIC}, l'identité
\code{4294967295} des sorties de la sonde. \code{analyze_cfg} construit un
\code{AnalysisCtx} par bloc (\cref{chap:cfg-analyzer-dominance}).

\rustsnippet{src/domains/context.rs}{127}{145}{un contexte sans requête ni inter-composants}

Le \code{static NULL} donne une référence \code{&'static dyn QueryContext}
sans allocation, \code{NullCtx} étant une structure unitaire. Le moteur
l'emploie avec des stores d'amorçage vides pour évaluer les constantes
de module, les arguments de hooks personnalisés et les initialiseurs de
\code{useState} (\srcline{src/engine/fixpoint.rs}{168},
\srcline{src/engine/fixpoint.rs}{245}, \srcline{src/engine/fixpoint.rs}{327}).

\subsection{\code{QueryContext} et ses deux implémentations}

\rustsnippet[label={lst:statevalue-stabilite-query}]{src/domains/context.rs}{27}{49}{le canal de requête entre domaines}

Le trait n'a plus qu'une requête, \code{callback_body} : le CFG du corps
d'un \code{useCallback}, que la réécriture en \code{CallbackVal(label)} a
sorti de l'arbre d'expressions et qui n'est donc pas atteignable par le tas
comme un \code{FnLit} ordinaire. Sans réponse (\code{None}), l'interpréteur
n'exécute pas le corps de l'appel
(\srcline{src/domains/interp/interpreter.rs}{579}). \code{NullCtx} répond
toujours \code{None} ; son commentaire \og returns \code{Top} for every
query\fg{} est daté. Pendant le point fixe, c'est \code{FixpointCtx} qui
répond :

\rustsnippet[label={lst:statevalue-stabilite-fixpointctx}]{src/domains/context.rs}{147}{171}{le contexte de requête du point fixe}

\begin{piege}
  \code{NullCtx} sert aussi aux trois amorçages du
  \cref{sec:statevalue-stabilite-contextes} qui précèdent la boucle de point
  fixe (constantes de module, retours d'arguments de hooks personnalisés,
  initialiseur paresseux d'un \code{useState}) : un appel à travers une
  variable liée par \code{useCallback}, rencontré pendant l'un de ces
  amorçages, n'y exécute donc pas son corps. Cela ne fuit dans aucun état
  retenu : le store et le tas passés à \code{AnalysisCtx::null} pour ces
  amorçages sont des copies jetables (\src{src/engine/fixpoint.rs}{164}{168},
  \src{src/engine/fixpoint.rs}{240}{244}, \src{src/engine/fixpoint.rs}{317}{320}),
  distinctes du store du point fixe et jamais fusionnées dedans. Ce qui reste
  ouvert est une question d'exécution de l'interpréteur, pas de domaine :
  est-ce qu'un effet de bord perdu de cette façon (un \code{useCallback} qui
  écrirait un \emph{autre} slot depuis un initialiseur paresseux) pourrait
  échapper au point fixe par un autre chemin ; le \cref{chap:transfert-stores}
  tranche.
\end{piege}

Seul \code{callbacks} est lu (le clonage est celui d'un \code{Arc}, en temps
constant) ; les champs \code{state} et \code{memo} ne sont lus par aucune
méthode du trait, et la phrase \og Local variables return \code{Bottom}\fg{}
décrit un comportement disparu.

\begin{decision}[ADR-007 --- requêtes entre domaines par un trait]
\textbf{Décide} (\og Decision A\fg{}, solution B3) : un trait
\code{QueryContext} passé en \code{&dyn} à chaque méthode de
\code{Transfer}, ce qui garde le trait utilisable en objet (object-safe).
\textbf{Écarte} : le \code{ask} générique de MOPSA, qui demande des GADT
absents de Rust stable (et la spécialisation, instable). \textbf{Prévoit}
(\og Decision B\fg{}, B1) : une migration vers un \og Manager\fg{} générique à
types marqueurs si le nombre de domaines dépassait cinq environ ; elle n'a
pas eu lieu. \textbf{Daté} : l'ADR décrit une requête \code{state_value_of}
et un troisième contexte \code{AnalysisQueryCtx} qui n'existent plus. L'ADR-015
a confirmé le rôle futur du trait : canal de lecture d'un éventuel produit
réduit au niveau de l'analyse.
\end{decision}

\subsection{\code{InterCtx} et \code{AnalyzeChildFn}}

Quand l'évaluateur rencontre un élément \code{<Child />} dont le composant
est connu, il analyse l'enfant avec les props de l'appelant
(\cref{chap:inter-composants}). Cela crée un cycle de modules : le transfert
(\code{domains}) doit appeler le point fixe (\code{engine}), qui appelle le
transfert. Le cycle est cassé par un pointeur de fonction fourni par le
moteur :

\rustsnippet{src/domains/context.rs}{15}{25}{le pointeur qui casse le cycle transfert--point fixe}

\rustsnippet[label={lst:statevalue-stabilite-interctx}]{src/domains/context.rs}{51}{103}{le contexte inter-composants}

Tout l'état partagé (cache des analyses d'enfants, store d'état partagé
entre composants, graphe d'appels, statistiques, résultats) passe par des
\code{RefCell}, pour que \code{InterCtx} circule en référence partagée sans
\code{&mut} imbriqués. \code{child} copie toutes les références et ne crée
qu'une pile neuve, où il empile le composant \emph{courant}.

\begin{invariant}{La pile d'appel contient les ancêtres}{statevalue-stabilite-callstack}
  La \code{call_stack} d'un \code{InterCtx} contient les ancêtres du
  composant courant, pas le composant lui-même : il n'y est poussé qu'en
  descendant, par \code{child}. D'où le second disjoint
  \code{self.component == id} de \code{is_recursive}.
\end{invariant}

L'évaluateur appelle \code{is_recursive} avant \code{child}
(\srcline{src/domains/transfer/state_value.rs}{521},
\srcline{src/domains/transfer/state_value.rs}{578}) : un composant récursif
n'est pas ré-analysé, son élément vaut \code{ref(Stable)} et la coupure est
comptée dans les statistiques. Les contextes racines sont construits par
\code{analyze_program} avec une pile vide
(\src{src/engine/fixpoint.rs}{729}{741}).

% ===========================================================================
\section{Limites et défauts observés}
\label{sec:statevalue-stabilite-limites}

Les défauts qui suivent ont été observés au commit \snapshotCommit{} sur
les fichiers de \file{/tmp/ch10/}. Ils sont présentés comme des constats ;
ceux qui n'ont pas d'issue sont signalés comme tels.

\subsection{Un mémo sur un état numérique : un FP de projection}

\begin{lstlisting}[language=TypeScript,caption={Deux mémos sur un état (\file{/tmp/ch10/memo.tsx}, extrait)},label={lst:statevalue-stabilite-memonum}]
export function MemoOverNumState() {
  const [n, setN] = useState(0);
  const doubled = useMemo(() => n * 2, [n]);
  useEffect(() => { console.log(doubled); }, [doubled]);
  return <button onClick={() => setN(5)}>{doubled}</button>;
}

export function MemoOverObjState() {
  const [o, setO] = useState({ a: 1 });
  const a = useMemo(() => ({ v: o.a }), [o]);
  useEffect(() => { console.log(a); }, [a]);
  return <button onClick={() => setO({ a: 2 })}>{a.v}</button>;
}
\end{lstlisting}

\begin{sortie}
  MemoOverNumState  (4 hooks)  /tmp/ch10/memo.tsx
    warn   always-unstable-deps  [hook:2]  (line 6:2)  this effect depends on `doubled`, a new reference every render, so `Object.is` always differs and the effect re-runs on every render regardless of the other deps
       → the value flows through `doubled`, bound here (line 5:8)
  MemoOverObjState  (4 hooks)  /tmp/ch10/memo.tsx  ✓
\end{sortie}

Le \Warning{} est faux : \code{doubled} est un nombre, recalculé seulement
quand \code{n} change. La sonde donne \code{state[0] = number[0, 5]} et
\code{memo[1] (Memo) = ref(PerRender)}. La chaîne causale se lit dans le code :
\begin{enumerate}
  \item la dep \code{n} n'est pas syntaxiquement \code{StateVal(0)} : le
    lowering lie \code{let n = StateVal(0)} (\srcline{src/lowering/hook_extractor.rs}{783})
    et la dep est la variable \code{n} ; la projection structurelle de
    \code{recompute_memo} ne s'applique donc pas ;
  \item la valeur de \code{n} est \code{number[0, 5]}, sans composante
    \code{reference} : la conversion côté lecture ne l'a pas étiquetée, et
    \code{versioned_reference} rend \code{None} ;
  \item \code{to_stability} projette l'intervalle non ponctuel sur
    \code{PerRender}, au sens \emph{may} de motion-wins ;
  \item \code{recompute_memo} range ce \code{PerRender} dans la composante
    \code{reference} du mémo, où \code{is_unstable_reference_only} le lit
    comme un fait \emph{must}.
\end{enumerate}
Le glissement de sens de \code{PerRender} passe ainsi d'une projection may à
un fait must. Avec un état objet, la lecture porte déjà \code{Versioned} dans
sa composante référence, et le problème ne se pose pas. Aucune issue
correspondante n'a été trouvée au commit \snapshotCommit{} ; la même cause
produit le FP de \cref{lst:statevalue-stabilite-memoprim}.

\subsection{Les cycles finis}

\begin{lstlisting}[language=TypeScript,caption={Un booléen qui bascule (\file{/tmp/ch10/kinds.tsx}, extrait)},label={lst:statevalue-stabilite-toggle}]
export function Toggle() {
  const [on, setOn] = useState(false);
  useEffect(() => {
    setOn(!on);
  }, [on]);
  return <div>{on ? 1 : 0}</div>;
}
\end{lstlisting}

C'est une vraie boucle infinie : chaque écriture change la valeur, la dep
change, l'effet se ré-exécute. L'analyseur déclare pourtant le composant
propre (la ligne \code{Toggle (2 hooks)} de la sortie de \code{--show-clean}
porte la coche des composants sans diagnostic). Le store converge sans widening sur \code{boolean}
($\Top$ de \code{BoolVal}) ; le bras point fixe d'\regle{infinite-loop} exige
un slot élargi, et le churn ne concerne que la sorte référence. Le produit
est un domaine d'\emph{ensembles de valeurs atteignables} : un cycle à
travers un nombre fini de valeurs a un point fixe abstrait fini, et
l'information \og la valeur change à chaque exécution\fg{}, une propriété de
\emph{transition}, n'y est pas représentable. Pour une référence, au
contraire, le domaine conserve un fait de transition (\code{PerRender},
\og neuve à chaque écriture\fg{}), ce qui permet l'\Error{} sur
\code{ObjChurn}. Le \cref{sec:treillis-simples-limites} décrit la même
famille pour \texttt{(i + 1) \% 3} et les chaînes ; ce FN n'est ni dans une
issue ni dans \file{docs/limitations.md}, et l'ADR-008 affirmait l'inverse.

\subsection{Une chaîne qui grandit : \code{is_unbounded} ignore \code{str}}

Le composant \code{StrGrow} de \file{/tmp/ch10/limits.tsx} écrit
\code{setS(s + "x")} dans un effet de deps \code{[s]}. Le store passe de
\code{{"a"}} à \code{{"a", "ax"}}, etc., jusqu'à dépasser quatre
constantes : \code{StrConst::Top} (sonde : \code{state[0] = string
widened=true}). Mais \code{is_unbounded}
(\cref{lst:statevalue-stabilite-unbounded}) ne regarde que \code{num}, la
référence et le résidu, et la garde n°2 du bras point fixe coupe :
\begin{sortie}
  StrGrow  (2 hooks)  /tmp/ch10/limits.tsx
    info   widening-info  (line 5:2)  state `s` kept changing during analysis and was approximated to converge, so findings that depend on it may be imprecise
    ...
    verified  infinite-loop  no effect diverges into an infinite render loop
\end{sortie}
(sortie de \code{--show-clean --trace --info}, élaguée). \og verified\fg{}
sur une vraie boucle infinie : un FN. Parmi les composantes qu'ignore
\code{is_unbounded}, \code{str} est la seule dont l'ensemble concret soit
infini, donc la seule qui puisse grandir sans borne. Aucune issue ni
mention dans \file{docs/limitations.md} n'a été trouvée.

\subsection{Un élément JSX natif classé \code{Stable} : un FN sans trace}
\label{sec:statevalue-stabilite-nativeelem}

La ligne \og élément JSX natif\fg{} du \cref{tab:statevalue-stabilite-sources}
range la même valeur que \code{useRef()} : \code{ref(Stable)}. C'est pourtant
la seule des huit lignes de ce tableau dont la colonne \og Raison\fg{} ne
renvoie à rien : ni fait must sur une allocation, ni garantie de React, ni
commentaire dans le code (\src{src/domains/transfer/state_value.rs}{153}{156}).
Or un élément JSX natif est une valeur allouée : \code{<div/>} désucre en un
appel (\code{React.createElement} ou \code{jsx}), qui rend un objet neuf à
chaque évaluation, exactement comme \code{{}} ou \code{() => {}} deux
lignes au-dessus du même tableau.

\begin{lstlisting}[language=TypeScript,caption={Un élément natif en dep (\file{/tmp/ch10/limits.tsx}, extrait)},label={lst:statevalue-stabilite-elemdep}]
export function ElemDep() {
  const el = <div />;
  useEffect(() => { console.log(el); }, [el]);
  return el;
}
\end{lstlisting}

\begin{sortie}
  ElemDep  (1 hooks)  /tmp/ch10/limits.tsx  ✓
\end{sortie}

Composant propre, alors que \code{el} est \code{Object.is}-différent à
chaque rendu et que l'effet devrait donc se ré-exécuter systématiquement.
Comparer avec le même programme où \code{el} vaut \code{{}} au lieu de
\code{<div/>} : \regle{always-unstable-deps} tire, sans surprise
(\cref{lst:statevalue-stabilite-objstate} donne le mécanisme). La cause est
locale à la ligne citée plus haut : \code{Expr::ObjectLit}, \code{ArrayLit}
et \code{FnLit} rendent \code{ref(PerRender)}, mais \code{Expr::NativeElem}
juste en dessous rend \code{ref(Stable)}, sans qu'aucune des trois autres
sources du produit (le tableau, un ADR, un test dédié) ne justifie l'écart.
\code{missing-deps} hérite du même trou : omettre \code{el} des deps ne
déclenche rien non plus, puisque \code{is_stable} le couvre.

\begin{piege}
  Une dep qui est elle-même un élément JSX natif est un motif rare en
  pratique (on ne compare pas des arbres React par identité dans une
  application normale), ce qui explique sans doute que le défaut n'ait
  jamais atteint un corpus de test. Cela n'en fait pas moins un faux
  négatif au sens strict de la \cref{def:soundness} : rien dans le code ne
  borne le domaine de \code{ElemDep} à des programmes \og improbables\fg{}.
  Aucune issue ne suit ce défaut au commit \snapshotCommit{}.
\end{piege}

\subsection{Limites connues et suivies}

\begin{itemize}
  \item \textbf{Slot jamais écrit} (\issue{41}, \code{precision-fp}) : un
    \code{useState({...})} dont aucun setter n'est appelé lit
    \code{Versioned}, pas \code{Stable}, et la dep ne peut pas être omise.
    Le raffinement demanderait un bit \og slot jamais écrit\fg{} après le
    point fixe ; l'ADR-020 (non-changement n°3) garde
    \code{may_written_slots} syntaxique, parce qu'un bit observé pourrait
    sous-compter sur un chemin élagué.
  \item \textbf{Constantes de module opaques} (\issue{34}) : \code{const X =
    f()} vaut $\Top$ ; le produit n'a pas d'encodage \og sorte inconnue,
    identité stable\fg{}.
  \item \textbf{Méthodes ambiguës} (\issue{22}) : \code{slice} et
    \code{concat} restent $\Top$ ; lever la limite demanderait la sorte du
    récepteur.
  \item \textbf{Écriture dérivée du slot écrit} (\issue{157}), vue au
    \cref{sec:statevalue-stabilite-lecture}.
  \item \textbf{Silence sur $\Top$} : \code{const v = flag ? { a: 2 } : obj} vaut \code{Unknown} (join de \code{PerRender} et de
    \code{Versioned}), et \regle{always-unstable-deps} se tait sur une dep
    \code{v} (composant \code{MixedRef} de \file{kinds.tsx}, propre). C'est
    le choix documenté de la règle (\og \code{Top} (precision lost) stays
    silent to avoid FPs\fg{},
    \src{src/rules/impls/always_unstable_deps.rs}{161}{163}) : elle
    n'affirme que des faits must.
\end{itemize}

\begin{piege}
  Plusieurs commentaires et ADR du périmètre sont datés, et il ne faut pas
  s'y fier sans relire le code : la documentation de \code{NullCtx}
  (\og returns \code{Top}\fg{}) et de \code{FixpointCtx} (\og Local variables
  return \code{Bottom}\fg{}) ; ADR-007 (\code{state_value_of},
  \code{AnalysisQueryCtx}) ; ADR-015 et ADR-017 (\code{Symbol} au lieu de
  \code{ComponentId}) ; ADR-002 (fichiers inexistants, seuil 2) ; ADR-008
  (oscillation booléenne détectée) ; et le nom de \code{from_init}.
\end{piege}

% ===========================================================================
\begin{resume}
\begin{itemize}
  \item Une valeur abstraite est un \textbf{produit ponctuel} sur les huit
    sortes disjointes de JavaScript (\code{StateValue},
    \file{src/domains/impls/state_value.rs}) : les sortes étant disjointes,
    la complétion disjonctive dégénère en produit, sans opérateur de
    réduction (ADR-015, qui remplace la somme plate d'ADR-008). Join, meet,
    widening et ordre sont composante par composante ; $\gam$ est l'union des
    concrétisations (\cref{def:statevalue}).
  \item \code{as_arith} applique \code{ToNumber(null) = 0}, ce qui fait
    détecter le compteur \code{useState(null)} ; \code{undefined}, les
    booléens et les chaînes rendent $\Top$ en arithmétique.
    \code{KindMask} et \code{populated_kinds} sont l'unique énumération des
    sortes.
  \item La composante \code{reference} porte une \textbf{stabilité},
    borne sur la \emph{trace de changement} d'une valeur au fil des rendus.
    Le treillis \code{Stability} d'ADR-017 (\cref{def:stabilite}) combine une
    borne may (\code{Stable} $\lleq$ \code{Versioned(S)} $\lleq$
    \code{VersionedTop}, $|S| \le 4$) et une borne must (\code{PerRender}),
    incomparables, jointes en \code{Unknown}. Son join est sound et de
    hauteur 8 ; son meet sous-estime l'intersection des $\gam$ et n'est pas
    utilisé.
  \item \textbf{Double vue} : le store range le join des valeurs écrites
    (\code{ref(PerRender)} pour un état objet) ; la lecture de
    \code{StateVal(l)} convertit la référence en \code{Versioned({l})}, en un
    seul point de l'évaluateur. C'est sound sous deux hypothèses couvertes
    ailleurs : pas de set pendant le rendu (\regle{setter-in-render}) et le
    bras churn d'\regle{infinite-loop}, qui certifie \code{ObjChurn}.
  \item Les littéraux, \code{new} et les allocateurs connus sont
    \code{PerRender} (must) ; \code{useRef} et le setter sont stables ; un
    mémo vaut le join des stabilités de ses deps (\code{recompute_memo}),
    rangé dans la composante \code{reference} : sa sorte réelle est perdue.
  \item Les règles lisent des \textbf{projections} :
    \code{is_unstable_reference_only} (must), \code{is_stable} (sûreté),
    \code{is_unbounded} (may), \code{is_finitely_valued},
    \code{to_stability} (motion-wins, où \code{PerRender} ne veut dire que
    \og peut changer à chaque rendu\fg{}) et \code{StabilityVerdict}.
  \item \code{AnalysisCtx}, \code{QueryContext}, \code{FixpointCtx},
    \code{InterCtx} et \code{AnalyzeChildFn} (\file{src/domains/context.rs})
    relient les domaines au moteur.
  \item Défauts observés : FP des mémos sur un état numérique, FN des cycles
    finis, des chaînes qui grandissent et de l'élément JSX natif classé
    \code{Stable} ; limites suivies \issue{41}, \issue{34}, \issue{22},
    \issue{157}.
\end{itemize}
Le \cref{chap:transfert-stores} décrit l'évaluateur complet qui fabrique ces
valeurs, et les stores (\cref{def:store}) qui les rangent.
\end{resume}

\section*{Exercices}
\addcontentsline{toc}{section}{Exercices}

\begin{exercice}{Calculer dans le produit}{statevalue-stabilite-produit}
  Soit $a = \code{number[0, 2]|null}$ et $b = \code{string{"a"}}$.
  (a) Calculer $a \join b$ et sa concrétisation. (b) Que valent
  \code{to_stability}, \code{is_unbounded} et \code{as_arith} sur
  $a \join b$ ? (c) Les valeurs $c = \code{number[0, 1]|null}$ et
  $d = \code{number[0, 2]}$ sont-elles comparables ?
\end{exercice}
\begin{solution}
  (a) Composante par composante : \code{number[0, 2]|string{"a"}|null}, de
  concrétisation $[0,2] \cup \{\code{"a"}\} \cup \{\code{null}\}$.
  (b) \code{to_stability} : l'intervalle n'est pas un point, donc le drapeau
  \og en mouvement\fg{} est levé ; trois sortes sont peuplées, ce qui joint
  \code{Unknown} ; motion-wins rend \code{PerRender}. \code{is_unbounded} :
  faux (bornes finies, pas de référence \code{PerRender}, pas de résidu).
  \code{as_arith} : \code{None} (la chaîne est peuplée), donc \code{n + 1}
  vaudrait $\Top$. (c) Non : sur \code{num}, $c \sqsubset d$ ; sur
  \code{null}, $c \sqsupset d$ ; \code{merge_ord} rend \code{None}.
\end{solution}

\begin{exercice}{Le meet trop petit}{statevalue-stabilite-meet}
  Exhiber un programme React dont une expression a une trace de changement
  à la fois dans $\gam(\code{Versioned}(\{l\}))$ et dans
  $\gam(\code{PerRender})$. Qu'en conclure sur \code{meet} ?
\end{exercice}
\begin{solution}
  La lecture \code{obj} dans \code{ObjChurn} (\cref{lst:statevalue-stabilite-objchurn}) :
  chaque rendu $r_i$ ($i \ge 1$) est précédé d'un \code{setObj} avec un objet
  neuf, donc $C = \{1,\dots,n\} = \mathit{Set}(\{\code{obj}\})$. La trace
  appartient aux deux concrétisations, alors que $\code{Versioned}(\{l\})
  \meet \code{PerRender} = \Bot$. Le meet du code est une borne inférieure
  dans l'ordre abstrait mais pas une sur-approximation de l'intersection des
  $\gam$ : s'en servir pour raffiner éliminerait ce comportement réel, un FN.
\end{solution}

\begin{exercice}{Casser l'invariant}{statevalue-stabilite-invariant}
  Un contributeur construit à la main \code{Stability::Versioned(S5)} avec
  $|S_5| = 5$. (a) Le join reste-t-il sound ? (b) Quelle propriété de la
  \cref{prop:statevalue-stabilite-lub} est perdue ?
\end{exercice}
\begin{solution}
  (a) Oui : la preuve du \cref{thm:statevalue-stabilite-join-stab} n'utilise
  pas la borne sur $|S|$ ; $\code{Versioned}(S_5)$ a simplement la
  concrétisation attendue. (b) La forme canonique : $\code{Versioned}(S_5)
  \join \code{Versioned}(S_5)$ rend $\code{Versioned}(S_5)$ par le bras
  d'égalité, alors que la même information s'écrirait canoniquement
  \code{VersionedTop}. La plus petite borne supérieure n'est plus garantie
  (un $\code{Versioned}(S_5)$ s'intercale entre les \code{Versioned} de
  taille 4 et \code{VersionedTop}), et la borne de hauteur 8 ne vaut plus
  pour les chaînes qui passent par des éléments non canoniques. La
  terminaison tient encore, parce que tout join de deux ensembles distincts
  repasse par \code{versioned}. C'est pourquoi les constructeurs \code{versioned} et \code{versioned_by}
  sont le seul chemin recommandé.
\end{solution}

\begin{exercice}{Projeter à la main}{statevalue-stabilite-projeter}
  Donner \code{to_stability} de : $[3,3]$ ;
  \code{number[3, 3]|null} ; \code{boolean} ($\Top$ de \code{BoolVal}) ;
  \code{ref(Versioned({a}))|undefined} ; \code{setter(...)} ; $\Top$ ;
  \code{number[0, inf]|ref(Stable)}.
\end{exercice}
\begin{solution}
  \code{Stable} (un point) ; \code{Unknown} (deux sortes) ; \code{Unknown}
  (booléen $\Top$) ; \code{Unknown} (le join partiel vaut
  \code{Versioned({a})}, mais deux sortes sont peuplées) ; \code{Stable}
  (identité garantie par React) ; \code{Unknown} (résidu) ; \code{PerRender}
  (motion-wins, l'intervalle n'est pas un point).
\end{solution}

\begin{exercice}{Corriger au bon niveau}{statevalue-stabilite-memonum}
  Expliquer pourquoi \code{MemoOverObjState} est propre et
  \code{MemoOverNumState} ne l'est pas
  (\cref{lst:statevalue-stabilite-memonum}). Pourquoi ajouter, dans
  \regle{always-unstable-deps}, un cas \og ignorer les deps qui sont des
  mémos\fg{} serait-il contraire aux principes du projet ? Où chercher une
  correction ?
\end{exercice}
\begin{solution}
  La lecture de \code{o} porte \code{Versioned({o})} dans sa composante
  référence, que \code{versioned_reference} transmet au mémo ; la lecture de
  \code{n} n'a pas de composante référence, donc aucune étiquette, et
  \code{to_stability} projette $[0,5]$ sur \code{PerRender} (sens may), que
  le mémo range comme un fait must. Un cas spécial dans la règle
  contournerait un défaut de la couche domaine (principe \og pas de
  workarounds\fg{}) et masquerait des vrais positifs (un mémo sur une dep
  réellement fraîche). La correction doit se placer là où le sens glisse :
  soit dans \code{recompute_memo}, qui ne devrait pas convertir une
  projection may en fait must, soit dans le domaine, qui n'a pas d'endroit
  où loger une étiquette de version pour une sorte non référence. Au commit
  \snapshotCommit{}, aucune issue ne suit ce défaut.
\end{solution}

\begin{exercice}{\code{is_unbounded} et les chaînes}{statevalue-stabilite-strgrow}
  Écrire une version de \code{is_unbounded} qui ferme le FN \code{StrGrow}.
  Change-t-elle quelque chose pour \code{StrToggle}
  (\code{setMode(mode === "a" ? "b" : "a")}) ? Et pour un état écrit par des
  handlers avec cinq constantes distinctes ?
\end{exercice}
\begin{solution}
  Ajouter le disjoint \code{self.str == StrConst::Top} (une chaîne de
  contenu inconnu peut grandir sans borne). \code{StrToggle} converge sur
  \code{string{"a", "b"}}, sous le seuil : inchangé (le FN des cycles finis
  demeure). Cinq constantes donnent aussi \code{StrConst::Top} : le prédicat
  deviendrait vrai alors que l'ensemble concret est fini. Qu'un diagnostic en
  résulte dépend des autres portes du bras point fixe
  (\cref{chap:infinite-loop}) ; le seuil de \code{StrConst} confond
  \og beaucoup de constantes\fg{} et \og croissance\fg{}, et c'est le prix
  de la hauteur finie.
\end{solution}

\begin{exercice}{Détecter \code{Toggle}}{statevalue-stabilite-toggle}
  Pourquoi aucun domaine d'\emph{ensembles de valeurs} ne peut-il détecter
  la boucle de \cref{lst:statevalue-stabilite-toggle} ? Quelle information
  faudrait-il ajouter ?
\end{exercice}
\begin{solution}
  L'ensemble des valeurs atteignables de \code{on} est $\{\code{false},
  \code{true}\}$, qui est aussi celui d'un composant qui écrit une seule
  fois \code{true} puis s'arrête ; les deux programmes ont le même point fixe
  abstrait. Ce qui les distingue est une propriété de \emph{transition} :
  \og la valeur écrite diffère toujours de la valeur lue\fg{}
  ($\code{!on} \ne \code{on}$). Il faudrait un fait relationnel entre la
  lecture et l'écriture d'un même slot, analogue pour les primitives à ce
  que \code{PerRender} fournit pour les références.
\end{solution}
